\chapter{Asymptotic Adventures}

\begin{center}\itshape
They must be true because, if they were not true, no one would have had the imagination to invent them.\\
\upshape --- G.~H.~Hardy, on formulas in Ramanujan's first letter, \emph{Ramanujan: Twelve Lectures} (1940), Lecture~I
\end{center}

In 1911 the young Srinivasa Ramanujan, still unknown outside a small circle of admirers in Madras, sent a question to the \emph{Journal of the Indian Mathematical
Society}. It became Question~294 of the third volume. Take the exponential series at $x=n$ and cut it in the middle:
\[
1+\frac n{1!}+\frac{n^2}{2!}+\dots+\frac{n^{n-1}}{(n-1)!}+\theta_n\frac{n^n}{n!}=\frac{e^n}2 .
\]
The first $n$ terms fall a little short of half the sum, and the claim was that the missing piece is always a fraction $\theta_n$
of the next term, with $\theta_n$ between $\frac13$ and $\frac12$. For $n=1$ this says $1+\theta_1=\frac e2$, so $\theta_1=0.359\ldots$ No proof came. It took seventeen
years before G\'abor Szeg\H o (1928) and George Neville Watson (1929) published proofs.

Thirty-three years later a graduate student called Donald Knuth was thinking about hash tables. Store keys in a table of $m$ cells;
each key gets a random address, and if that cell is taken it goes to the next free cell, wrapping around at the end. How many cells
does a search inspect on average? The answer is written with sums like
\[
Q(n)=1+\frac{n-1}n+\frac{(n-1)(n-2)}{n^2}+\cdots,
\]
and Knuth needed their size. ``Linear probing was the first algorithm that I was able to analyze successfully'', he wrote in
1998, recalling how satisfying it had been when he first studied the problem in 1962. The function $Q$ turned out to be Ramanujan's: the cut exponential series is $\frac{n^n}{n!}Q(n)$. You
will prove that $\theta_n\to\frac13$ in Problem~\ref{pb:ramanujan}, and see where the unexpected $\frac13$ comes from.

This chapter is about such questions: not the exact value of a number, which is often hopeless, but its size when a parameter is
large. The tools come from Tome~0 (Stirling's formula, iterated sequences), Tome~3 (birthdays, coupons, local limit theorems) and
Tome~4, Chapter~11 (Laplace's method, saddle points, Euler--Maclaurin, the transfer theorem). We count far away with saddle
points and poles, wait for coincidences, iterate, and finally let chance run for a long time.

\begin{think}
For $n=1$, Ramanujan's $\theta_1$ is $0.359$; a computer gives $\theta_{10}=0.336$. Guess the limit, and guess how fast it is reached. People
enter a room one at a time. On average $24.6$ of them arrive before two share a birthday (Tome~3, Exercise~2.7). How many arrive,
on average, before \emph{three} share a birthday: about $37$, about $50$, or about $90$?
\end{think}

Throughout, $f\sim g$ means $f/g\to1$, $[z^n]F$ is the coefficient of $z^n$ in a power series $F$, and $n^{\underline k}=n(n-1)\cdots(n-k+1)$.
Unless stated otherwise, limits are as $n\to\infty$.

\section{Counting from far away}

Tome~4, Chapter~11 read the size of a coefficient from its generating function, by a saddle point or by the nearest singularity.
The first problems use both ideas, together with plain Laplace sums, on numbers from combinatorics.

\begin{exercise}[$\star$]\label{pb:powersum}
\emph{The last terms win.} Show that
\[
\frac1{n^n}\sum_{k=1}^nk^n=\frac e{e-1}-\frac{e(e+1)}{2(e-1)^3}\cdot\frac1n+O\Bigl(\frac1{n^2}\Bigr).
\]
(Write the sum as $\sum_{j=0}^{n-1}(1-\frac jn)^n$.) For $n=100$ the left side is $1.57211$ and the two-term approximation is $1.57202$.
\end{exercise}

\begin{exercise}[$\star\star\star$]\label{pb:involasym}
\emph{How many involutions?} Let $t_n$ be the number of involutions of $\{1,\dots,n\}$, so that $\sum_nt_n\frac{z^n}{n!}=T(z)=e^{z+z^2/2}$ (Tome~4,
Exercise~2.4). Let $r=r_n>0$ be the root of $r+r^2=n$, and $b=r+2r^2$.
\begin{enumerate}[label=(\alph*)]
\item Show that $\frac{t_n}{n!}=\frac1{2\pi}\int_{-\pi}^\pi T(re^{i\theta})\,r^{-n}e^{-in\theta}\dd\theta$.
\item Show that $|T(re^{i\theta})|\le T(r)\exp\bigl(-2r\sin^2\frac\theta2-r^2\sin^2\theta\bigr)$.
\item Show that
\[
\frac{t_n}{n!}\sim\frac{T(r)}{r^n\sqrt{2\pi b}}.
\]
(Cut the integral at $|\theta|=r^{-4/5}$.)
\item Deduce that $t_n\sim\frac1{\sqrt2}\,n^{n/2}e^{-n/2+\sqrt n-1/4}$.
\end{enumerate}
\end{exercise}

\begin{exercise}[$\star\star\star$]\label{pb:bellasym}
\emph{The size of the Bell numbers.} Let $b_n$ be the number of partitions of an $n$-set, so that $b_n=\frac1e\sum_{k\ge0}\frac{k^n}{k!}$ (Dobi\'nski's
formula, Tome~4, Chapter~2). For $n\ge3$ let $u=u_n$ be the positive root of $ue^u=n$.
\begin{enumerate}[label=(\alph*)]
\item Show that $k^n\le\bigl(\frac n{eu}\bigr)^ne^{uk}$ for every integer $k\ge0$, and deduce that $\ln b_n\le n\bigl(u-1+\frac1u\bigr)-1$.
\item Using the single term $k=\lceil e^u\rceil$ of Dobi\'nski's formula, show that $\ln b_n\ge n\bigl(u-1+\frac1u\bigr)-O(\ln n)$.
\item With $L=\ln n$, show that $u=L-\ln L+\frac{\ln L}L+O\bigl((\frac{\ln L}L)^2\bigr)$, and deduce that
\[
\ln b_n=n\Bigl(\ln n-\ln\ln n-1+\frac{1+\ln\ln n}{\ln n}+O\Bigl(\Bigl(\frac{\ln\ln n}{\ln n}\Bigr)^2\Bigr)\Bigr).
\]
\end{enumerate}
\end{exercise}

\begin{exercise}[$\star\star\star$]\label{pb:distinctcycles}
\emph{All cycles of different lengths.} Let $d_n$ be the probability that a uniform random permutation of $\{1,\dots,n\}$ has cycles of pairwise distinct lengths.
\begin{enumerate}[label=(\alph*)]
\item Show that $\sum_{n\ge0}d_nz^n=\prod_{k\ge1}\bigl(1+\frac{z^k}k\bigr)$ for $|z|<1$, and compute $d_1,\dots,d_5$.
\item Show that this product equals $\frac{h(z)}{1-z}$, where $h(z)=\prod_{k\ge1}\bigl(1+\frac{z^k}k\bigr)e^{-z^k/k}$, and that $h(1)=e^{-\gamma}$.
\item For a power series $f=\sum a_jz^j$ write $\|f\|=\sum_j|a_j|$. Show that $\|fg\|\le\|f\|\,\|g\|$ and that $\bigl\|(1+\frac{z^k}k)e^{-z^k/k}-1\bigr\|\le\frac1{k^2}$.
Deduce that the coefficients $h_j$ of $h$ satisfy $\sum_j|h_j|<\infty$ and $\sum_jh_j=h(1)$.
\item Deduce that $d_n\to e^{-\gamma}=0.5614\ldots$
\end{enumerate}
\end{exercise}

\begin{exercise}[$\star\star\star$]\label{pb:zigzag}
\emph{Up, down, up, down.} A permutation $\sigma$ of $\{1,\dots,n\}$ is \emph{alternating} if $\sigma(1)<\sigma(2)>\sigma(3)<\sigma(4)>\cdots$. Let $A_n$ be the number of
alternating permutations, with $A_0=A_1=1$; the values continue $1,2,5,16,61,272,1385$.
\begin{enumerate}[label=(\alph*)]
\item Show that $2A_{n+1}=\sum_{k=0}^n\binom nkA_kA_{n-k}$ for $n\ge1$. (Look at where $n+1$ sits, and count the permutations that zigzag starting either up or down.)
\item Let $f(z)=\sum_nA_n\frac{z^n}{n!}$. Show that $2f'=1+f^2$ and $f(0)=1$, and deduce that $f(z)=\sec z+\tan z$ for $|z|<\frac\pi2$.
\item Show that $f(z)-\frac4{\pi-2z}$ extends to a holomorphic function on the disc $|z|<\frac{3\pi}2$, and deduce that
\[
\frac{A_n}{n!}=2\Bigl(\frac2\pi\Bigr)^{n+1}\bigl(1+O(\delta^n)\bigr)\qq\text{for every }\delta>\tfrac13.
\]
\item What is the probability that a random permutation of $10$ elements is alternating? Compare with the approximation.
\end{enumerate}
\end{exercise}

\begin{exercise}[$\star\star$]\label{pb:binompowers}
\emph{Powers of binomial coefficients.} Let $s\ge2$ be an integer. Show that
\[
\sum_{k=0}^n\binom nk^s\sim\frac{2^{sn}}{\sqrt s}\Bigl(\frac2{\pi n}\Bigr)^{(s-1)/2}.
\]
(Use the local de Moivre--Laplace theorem, Tome~3, Theorem~6.12, near the middle, and a cruder bound away from it.) Check the case $s=2$
against $\binom{2n}n$, and write the cases $s=3$ and $s=4$ explicitly.
\end{exercise}

\section{Waiting for coincidences}

Draw repeatedly, uniformly and independently, from $n$ values: birthdays, hash addresses, coupons. How long until something happens?
One trick runs through this section: let the draws arrive at random times, so that the $n$ counts become independent.

\begin{exercise}[$\star\star$]\label{pb:poissonise}
\emph{Three people with the same birthday.} Draw uniformly and independently from $n$ values, and let $T$ be the number of draws needed
until some value has appeared $m$ times ($m\ge2$). Let $e_m(x)=\sum_{j<m}\frac{x^j}{j!}$.
\begin{enumerate}[label=(\alph*)]
\item Show that $P(T>k)=\frac{k!}{n^k}[x^k]\,e_m(x)^n$ for every $k\ge0$.
\item Deduce that $\mathbb E[T]=\int_0^\infty e^{-t}e_m(t/n)^n\dd t$.
\item For $m=2$, deduce that $\mathbb E[T]=1+Q(n)$, where $Q(n)=\sum_{k=1}^n\frac{n^{\underline k}}{n^k}$, and recover Tome~3, Exercise~2.7.
\item For $m=3$, show that $\mathbb E[T]\sim\Gamma(\frac43)(6n^2)^{1/3}$.
\item For $n=365$ the exact values are $24.62$ for a pair and $88.74$ for a triple. What does (d) give? Answer the think box.
\end{enumerate}
\end{exercise}

\begin{exercise}[$\star\star\star$]\label{pb:ramanujan}
\emph{Ramanujan's $\frac13$.} Let $Q(n)$ be as in Problem~\ref{pb:poissonise}, let $R(n)=1+\frac n{n+1}+\frac{n^2}{(n+1)(n+2)}+\cdots$, and define $\theta_n$ by
$\sum_{k=0}^{n-1}\frac{n^k}{k!}+\theta_n\frac{n^n}{n!}=\frac{e^n}2$.
\begin{enumerate}[label=(\alph*)]
\item Show that $\frac{n^n}{n!}Q(n)=\sum_{k<n}\frac{n^k}{k!}$ and $\frac{n^n}{n!}R(n)=\sum_{k\ge n}\frac{n^k}{k!}$. Deduce that $Q(n)+R(n)=\frac{n!\,e^n}{n^n}$ and $\theta_n=\frac{R(n)-Q(n)}2$.
\item Show that
\[
Q(n)=n\int_1^\infty s^{n-1}e^{n(1-s)}\dd s,\qq R(n)=n\int_0^1s^{n-1}e^{n(1-s)}\dd s.
\]
\item Let $\varphi(s)=s-1-\ln s$ and $v(s)=\operatorname{sign}(s-1)\sqrt{2\varphi(s)}$, an increasing bijection from $\left]0,\infty\right[$ onto $\R$; write $s(v)$ for its inverse.
Show that
\[
Q(n)=n\int_0^\infty e^{-nv^2/2}g(v)\dd v,\qq R(n)=n\int_{-\infty}^0e^{-nv^2/2}g(v)\dd v,\qq g(v)=\frac v{s(v)-1},
\]
and that $0<g(v)\le1+|v|$ and $g(v)=1-\frac v3+O(v^2)$, uniformly on $\R$.
\item Deduce that $Q(n)=\sqrt{\frac{\pi n}2}-\frac13+O\bigl(\frac1{\sqrt n}\bigr)$ and $R(n)=\sqrt{\frac{\pi n}2}+\frac13+O\bigl(\frac1{\sqrt n}\bigr)$, so that $\theta_n\to\frac13$.
\item Deduce that $e^{-n}\sum_{k=0}^n\frac{n^k}{k!}=\frac12+\frac2{3\sqrt{2\pi n}}+O\bigl(\frac1n\bigr)$.
\end{enumerate}
\end{exercise}

Figure~\ref{fig:theta} shows how fast Ramanujan's numbers settle.

\begin{figure}[ht]
\centering
\fig{ch12-theta}
\caption{Ramanujan's $\theta_n$ for $n=1,\dots,40$ (dots), decreasing towards $\frac13$. The dashed curve is $\frac13+\frac4{135n}$, the next term of the expansion,
which the notebook checks numerically (Problem~\ref{pb:ramanujan}).}
\label{fig:theta}
\end{figure}

\begin{exercise}[$\star\star\star$]\label{pb:dixie}
\emph{Collecting every coupon twice.} Coupons of $n$ types are drawn uniformly and independently, and $T_m$ is the number of draws until
every type has been seen at least $m$ times. Let $p_m(x)=e^{-x}e_m(x)$, the probability that a Poisson variable of mean $x$ is less than $m$.
\begin{enumerate}[label=(\alph*)]
\item Let $N_t$ be a Poisson variable of mean $t$, independent of the draws. Show that $\mathbb E[T_m]=\int_0^\infty P(T_m>N_t)\dd t$, and deduce that
\[
\mathbb E[T_m]=n\int_0^\infty\bigl(1-(1-p_m(x))^n\bigr)\dd x.
\]
\item For $m=1$, recover $\mathbb E[T_1]=nH_n$ (Tome~3, Theorem~2.14).
\item For fixed $m$, show that $\mathbb E[T_m]=n\ln n+(m-1)\,n\ln\ln n+O(n)$.
\end{enumerate}
\end{exercise}

\begin{exercise}[$\star\star$]\label{pb:mappings}
\emph{Random maps.} Let $f$ be a uniform random map from $\{1,\dots,n\}$ to itself. Its \emph{functional graph} has an arrow $x\to f(x)$ for each $x$;
each connected component contains exactly one cycle, and $x$ is \emph{cyclic} if $f^k(x)=x$ for some $k\ge1$.
\begin{enumerate}[label=(\alph*)]
\item Show that $x$ lies on a cycle of length exactly $k$ with probability $\frac{n^{\underline k}}{n^{k+1}}$. Deduce that the expected number of cyclic points is
$Q(n)\sim\sqrt{\pi n/2}$.
\item Show that the expected number of components is $C_n=\sum_{k=1}^n\frac{n^{\underline k}}{k\,n^k}$.
\item Show that $C_n=\frac12\bigl(\ln(2n)+\gamma\bigr)+o(1)$. (For $n=100$, $C_n=2.979$ and the approximation is $2.938$.)
\end{enumerate}
\end{exercise}

\begin{exercise}[$\star\star$]\label{pb:parking}
\emph{Parking on a one-way street.} Spots $1,\dots,m$ line a one-way street. Cars arrive one by one; car $i$ prefers spot $a_i$, drives to it, and
parks in the first free spot from $a_i$ on, leaving the street if there is none. A \emph{parking function} of length $n$ is a sequence
$(a_1,\dots,a_n)\in\{1,\dots,n\}^n$ with which, for $m=n$, every car parks.
\begin{enumerate}[label=(\alph*)]
\item Put $n+1$ spots on a circle instead, with preferences in $\Z/(n+1)\Z$, so that every car parks and one spot stays empty. By rotating preferences,
show that there are $(n+1)^{n-1}$ parking functions of length $n$.
\item Back on the street, take $m=n+1$ and preferences uniform in $\{1,\dots,n+1\}^n$, and let $K+1$ be the first spot left empty. Show that
\[
P(K=k)=\binom nk\frac{(k+1)^{k-1}(n-k)^{n-k}}{(n+1)^n},
\]
and deduce Abel's identity $\sum_{k=0}^n\binom nk(k+1)^{k-1}(n-k)^{n-k}=(n+1)^n$.
\item Show that $P(K=k)\to p_k=\frac{(k+1)^{k-1}}{k!}e^{-(k+1)}$ for each fixed $k$, and that $\sum_kp_k=1$. (Use the tree function $T(z)=\sum_{m\ge1}m^{m-1}\frac{z^m}{m!}$,
which satisfies $T(z)e^{-T(z)}=z$.)
\item Show that $p_k\sim\frac1{\sqrt{2\pi}}k^{-3/2}$. So $K$ is usually small, yet the limit law has infinite mean.
\end{enumerate}
\end{exercise}

\section{Iterates, roots and constants}

Some sequences have no formula at all: they are defined by iterating a map, or as roots of equations. Their asymptotics come from
bootstrapping: a crude estimate, fed back into the equation, gives a better one. Others have a formula but hide a constant, and
Euler--Maclaurin names it.

\begin{exercise}[$\star\star$]\label{pb:xplusinv}
\emph{Adding the reciprocal.} Let $x_0=1$ and $x_{n+1}=x_n+\frac1{x_n}$.
\begin{enumerate}[label=(\alph*)]
\item Show that $x_n^2\ge2n+1$, and that $x_n^2\le2n+2+\ln n$ for $n\ge1$.
\item Show that $x_n^2=2n+\frac12\ln n+C+o(1)$ for a constant $C$. (The notebook finds $C=1.72314\ldots$)
\item Deduce that $x_n=\sqrt{2n}+\frac{\ln n+2C}{4\sqrt{2n}}+o\bigl(\frac1{\sqrt n}\bigr)$.
\end{enumerate}
\end{exercise}

\begin{exercise}[$\star\star\star$]\label{pb:sinesecond}
\emph{The sine iteration, second term.} Tome~0, Exercise~5.4 showed that $u_{n+1}=\sin u_n$, with $0<u_0\le\frac\pi2$, satisfies $u_n\sim\sqrt{3/n}$.
\begin{enumerate}[label=(\alph*)]
\item Show that $\frac1{\sin^2u}=\frac1{u^2}+\frac13+\frac{u^2}{15}+O(u^4)$ as $u\to0$.
\item Show that $\frac1{u_n^2}=\frac n3+O(\ln n)$, then that $\frac1{u_n^2}=\frac n3+\frac15\ln n+C+o(1)$ for a constant $C$ depending on $u_0$.
\item Deduce that $u_n=\sqrt{\frac3n}-\frac{3\sqrt3}{10}\cdot\frac{\ln n}{n^{3/2}}+O\bigl(\frac1{n^{3/2}}\bigr)$.
\end{enumerate}
\end{exercise}

\begin{exercise}[$\star$]\label{pb:tanroots}
\emph{The roots of $\tan x=x$.}
\begin{enumerate}[label=(\alph*)]
\item Show that for each $k\ge1$ the equation $\tan x=x$ has exactly one root $x_k$ in $\left]k\pi-\frac\pi2,k\pi+\frac\pi2\right[$, and that $k\pi<x_k<k\pi+\frac\pi2$.
\item With $a=(k+\frac12)\pi$, show that $x_k=a-\frac1a-\frac2{3a^3}+O(a^{-5})$.
\item The notebook finds $x_1=4.49341$ against $4.49381$ from (b). Explain why the approximation is already good for $k=1$.
\end{enumerate}
\end{exercise}

\begin{exercise}[$\star\star$]\label{pb:nbonacci}
\emph{Fibonacci with $n$ terms.} For $n\ge2$ let $P_n(x)=x^n-x^{n-1}-\dots-x-1$.
\begin{enumerate}[label=(\alph*)]
\item Show that $P_n$ has exactly one positive root $r_n$, that $1<r_n<2$, and that $r_n<r_{n+1}$. What is $r_2$?
\item Show that $r_n=2-2^{-n}-n\,2^{-2n-1}+O(n^22^{-3n})$. (Use $(x-1)P_n(x)=x^{n+1}-2x^n+1$.)
\item The number of coin-toss sequences of length $\ell$ without $n$ consecutive heads satisfies $a_\ell=a_{\ell-1}+\dots+a_{\ell-n}$. Explain why it grows like
$r_n^\ell$, and why $r_n$ is so close to $2$.
\end{enumerate}
\end{exercise}

\begin{exercise}[$\star\star$]\label{pb:sqrtsum}
\emph{The constant in $1+\sqrt2+\dots+\sqrt n$.} Use the continuation of $\zeta$ by Euler--Maclaurin (Tome~4, Proposition~11.7) to show that
\[
\begin{aligned}\sum_{k=1}^n\sqrt k&=\frac23n^{3/2}+\frac12\sqrt n+\zeta\Bigl(-\frac12\Bigr)+\frac1{24\sqrt n}+O\bigl(n^{-5/2}\bigr),\\
\sum_{k=1}^n\frac1{\sqrt k}&=2\sqrt n+\zeta\Bigl(\frac12\Bigr)+\frac1{2\sqrt n}-\frac1{24n^{3/2}}+O\bigl(n^{-7/2}\bigr),\end{aligned}
\]
where $\zeta(-\frac12)=-0.2078862\ldots$ and $\zeta(\frac12)=-1.4603545\ldots$ So the constant left over when the growing terms are removed is a value of the zeta function.
\end{exercise}

\begin{exercise}[$\star\star\star$]\label{pb:hyperfactorial}
\emph{The hyperfactorial and Glaisher's constant.} Let $H(n)=1^1\,2^2\,3^3\cdots n^n$.
\begin{enumerate}[label=(\alph*)]
\item Differentiate the identity of Tome~4, Proposition~11.7 with respect to $s$ and set $s=-1$, to show that
\[
\ln H(n)=\Bigl(\frac{n^2}2+\frac n2+\frac1{12}\Bigr)\ln n-\frac{n^2}4+\ln A+\frac1{720n^2}+O\Bigl(\frac1{n^4}\Bigr),
\]
where $\ln A=\frac1{12}-\zeta'(-1)$.
\item The notebook computes $A=1.2824271291\ldots$, Glaisher's constant. Check (a) numerically for $n=10$.
\end{enumerate}
\end{exercise}

\section{Chance in the long run}

The last problems let a random process run for a long time and ask what survives: a probability, a proportion, a growth rate.
Sometimes a formula exists; sometimes only a limit does, and the problem is to prove that it exists and to bound it.

\begin{exercise}[$\star\star$]\label{pb:diehit}
\emph{Will the running total hit $n$?} Roll a fair die repeatedly and add the results. Let $u_n$ be the probability that some running total equals $n$, with $u_0=1$.
\begin{enumerate}[label=(\alph*)]
\item Show that $u_n=\frac16(u_{n-1}+\dots+u_{n-6})$ for $n\ge1$ (with $u_j=0$ for $j<0$), and that $\sum_nu_nz^n=\frac1{1-D(z)}$, where $D(z)=\frac16(z+z^2+\dots+z^6)$.
\item Show that $1-D$ has a simple root at $z=1$ and no other root in the closed unit disc.
\item Deduce that $u_n\to\frac27$, the inverse of the mean roll $\frac72$, and that $u_n-\frac27=O(\rho^n)$ for some $\rho<1$. (The notebook finds $\rho=0.7302$.)
\item Prove the same for any law on $\{1,\dots,N\}$ with mean $\mu$ whose support has greatest common divisor $1$: then $u_n\to\frac1\mu$. What goes wrong if the gcd is $2$?
\end{enumerate}
\end{exercise}

\begin{figure}[ht]
\centering
\fig{ch12-diehit}
\caption{The probability $u_n$ that the running total of a die hits $n$, for $n=1,\dots,30$. It climbs as $\frac16(\frac76)^{n-1}$ up to $n=6$, then oscillates and
settles at $\frac27$ (dashed) with error $O(0.731^n)$ (Problem~\ref{pb:diehit}).}
\label{fig:diehit}
\end{figure}

\begin{exercise}[$\star\star$]\label{pb:banach}
\emph{Banach's matchboxes.} A smoker carries two boxes of $N$ matches, one in each pocket, and each time takes a match from a pocket chosen at random. At the moment he first
finds a box empty, let $R$ be the number of matches in the other.
\begin{enumerate}[label=(\alph*)]
\item Show that $P(R=r)=\binom{2N-r}N2^{-(2N-r)}$ for $0\le r\le N$.
\item Writing $a_r=P(R=r)$, show that $(2N-r)a_{r+1}=2(N-r)a_r$, and deduce that $\mathbb E[R]=(2N+1)\binom{2N}N4^{-N}-1$.
\item Deduce that $\mathbb E[R]=2\sqrt{N/\pi}-1+O\bigl(\frac1{\sqrt N}\bigr)$. With $50$ matches per box, how many are left on average?
\end{enumerate}
\end{exercise}

\begin{exercise}[$\star\star$]\label{pb:flory}
\emph{Dimers falling on a line.} A row has $n$ sites, numbered $1$ to $n$. Give each of the $n-1$ pairs $\{i,i+1\}$ an independent uniform label in $[0,1]$, and go
through the pairs in increasing order of labels, filling both sites of a pair whenever both are still empty. (At each step the pair filled is
uniform among the pairs still available.) Let $V_n$ be the expected number of sites left empty.
\begin{enumerate}[label=(\alph*)]
\item Show that $V_0=0$, $V_1=1$ and $(n-1)V_n=2\sum_{j=0}^{n-2}V_j$ for $n\ge2$.
\item Show that $\sum_{n\ge0}V_nx^n=\frac{xe^{-2x}}{(1-x)^2}$, so that $V_n=\sum_{k=0}^{n-1}(n-k)\frac{(-2)^k}{k!}$.
\item Deduce that $\bigl|V_n-(n+2)e^{-2}\bigr|\le\frac{2^{n+1}}{(n+1)!}$. In a long row, a fraction $e^{-2}=13.5\%$ of the sites stays empty.
\item Check the first-order result against the transfer theorem of Tome~4, Theorem~11.13.
\end{enumerate}
\end{exercise}

\begin{exercise}[$\star\star$]\label{pb:saw}
\emph{Self-avoiding walks.} Let $c_n$ be the number of $n$-step walks on $\Z^2$ from the origin, with unit steps, that never visit a point twice. So $c_1=4$, $c_2=12$, $c_3=36$, $c_4=100$.
\begin{enumerate}[label=(\alph*)]
\item Show that $c_{m+n}\le c_mc_n$.
\item (Fekete's lemma.) Show that if a real sequence satisfies $a_{m+n}\le a_m+a_n$, then $\frac{a_n}n\to\inf_{k\ge1}\frac{a_k}k$ (which may be $-\infty$).
\item Deduce that $\mu=\lim c_n^{1/n}$ exists and equals $\inf_nc_n^{1/n}$.
\item Show that $1+\sqrt2\le\mu\le3$. (For the lower bound, count the walks that use only the steps east, north and south and never follow north by south or south by north.)
\item The notebook counts $c_{14}=2374444$. What upper bound for $\mu$ does that give?
\end{enumerate}
\end{exercise}

\begin{figure}[ht]
\centering
\fig{ch12-saw}
\caption{Self-avoiding walks on $\Z^2$: $c_n^{1/n}$ (blue) and $\frac{c_{n+1}}{c_n}$ (orange) for $n\le14$, with the bounds $1+\sqrt2$ and $3$ of Problem~\ref{pb:saw} and the
numerical value $\mu\approx2.638$ (dashed). Since $\mu=\inf_nc_n^{1/n}$, the blue curve always stays above $\mu$; it is still far from it at $n=14$.}
\label{fig:saw}
\end{figure}

\begin{exercise}[$\star\star\star$]\label{pb:lcs}
\emph{Common subsequences of random words.} Let $A$ and $B$ be independent uniform words of length $n$ over $\{0,1\}$, and $L_n$ the length of their longest common
subsequence (not necessarily contiguous).
\begin{enumerate}[label=(\alph*)]
\item Show that $\mathbb E[L_{m+n}]\ge\mathbb E[L_m]+\mathbb E[L_n]$, and deduce that $\frac{\mathbb E[L_n]}n$ converges to $\gamma_2=\sup_n\frac{\mathbb E[L_n]}n$.
\item By exact enumeration the notebook finds $\mathbb E[L_8]=\frac{44803}{8192}$. What lower bound for $\gamma_2$ follows?
\item Show that $P(L_n\ge k)\le\binom nk^22^{-k}$. Let $H(\kappa)=-\kappa\ln\kappa-(1-\kappa)\ln(1-\kappa)$; using $\binom nk\le e^{nH(k/n)}$, show that $\gamma_2\le\kappa^*$, where $\kappa^*$ is the root in
$\left]\frac12,1\right[$ of $2H(\kappa)=\kappa\ln2$. Compute $\kappa^*$.
\end{enumerate}
\end{exercise}

\section{Hints}

\begin{hint}{pb:powersum}
For $0\le x\le\frac12$, $-x-\frac{x^2}2-x^3\le\ln(1-x)\le-x-\frac{x^2}2$. Treat $j\le\frac n2$ with this, and bound the terms with $j>\frac n2$ by $2^{-n}$. You will need
$\sum_jj^2x^j=\frac{x(1+x)}{(1-x)^3}$.
\end{hint}

\begin{hint}{pb:involasym}
(a) Cauchy's formula on the circle $|z|=r$. (b) Compute $\operatorname{Re}(re^{i\theta}+\frac12r^2e^{2i\theta})$. (c) Expand $re^{i\theta}+\frac12r^2e^{2i\theta}-in\theta$ to order $2$ in $\theta$: the
linear term vanishes because $r+r^2=n$. (d) $r=\sqrt n-\frac12+\frac1{8\sqrt n}+O(n^{-3/2})$, $r+\frac{r^2}2=\frac{n+r}2$, and Stirling for $n!$.
\end{hint}

\begin{hint}{pb:bellasym}
(a) Maximise $x^ne^{-ux}$ over $x\ge0$, then sum $\frac{e^{uk}}{k!}$. Remember $\ln n-\ln u=u$ and $e^u=\frac nu$. (b) Use $\ln m!\le(m+\frac12)\ln m-m+1$ (Tome~0), and note that $x\ln x-x$
increases by at most $\ln(x+1)$ when $x$ increases by at most $1$. (c) $u=L-\ln u$ and $L-\ln L\le u\le L$.
\end{hint}

\begin{hint}{pb:distinctcycles}
(a) The exponential formula (Tome~4, Theorem~2.9), each cycle length being used at most once. (b) $\prod_ke^{z^k/k}=\frac1{1-z}$, and $\prod_{k\le N}(1+\frac1k)=N+1$.
(c) $(1+w)e^{-w}=1-\sum_{j\ge2}(-1)^j\frac{j-1}{j!}w^j$. The partial products form a Cauchy sequence for $\|\cdot\|$. (d) $d_n=h_0+\dots+h_n$.
\end{hint}

\begin{hint}{pb:zigzag}
(a) Removing $n+1$ from a permutation that zigzags (up first or down first) leaves two blocks; read the left one backwards. (b) Both sides of $2f'=1+f^2$ have the
same coefficients by (a); check that $\sec+\tan=\frac{1+\sin z}{\cos z}$ solves the equation, and compare coefficients. (c) The zeros of $\cos$ nearest to $0$ are $\pm\frac\pi2$; at
$-\frac\pi2$ the numerator $1+\sin z$ vanishes too. Use Cauchy's estimates (Tome~4, Corollary~9.12).
\end{hint}

\begin{hint}{pb:binompowers}
With $p_k=\binom nk2^{-n}$, the sum is $2^{sn}\sum_kp_k^s$. Near the middle, $p_k\approx\frac2{\sqrt n}\phi(x_k)$ with $x_k=\frac{2k-n}{\sqrt n}$, and the $x_k$ are spaced $\frac2{\sqrt n}$ apart.
Away from it, $\sum p_k^s\le(\max_kp_k)^{s-1}P(|x_K|>A)$. Finally $\int_\R\phi^s=(2\pi)^{(1-s)/2}s^{-1/2}$.
\end{hint}

\begin{hint}{pb:poissonise}
(a) Exponential generating functions count words. (b) $k!=\int_0^\infty t^ke^{-t}\dd t$, and the sum over $k$ is finite. (c) Expand $(1+\frac tn)^n$. (d) Substitute $t=nx$; the
function $e^{-x}(1+x+\frac{x^2}2)$ is $1-\frac{x^3}6+O(x^4)$ near $0$.
\end{hint}

\begin{hint}{pb:ramanujan}
(a) $\frac{n^{n-j}}{(n-j)!}=\frac{n^n}{n!}\cdot\frac{n^{\underline j}}{n^j}$. (b) $\sum_{k\ge n}e^{-x}\frac{x^k}{k!}=\int_0^x\frac{t^{n-1}e^{-t}}{(n-1)!}\dd t$ (differentiate both sides), then put $x=n$ and $t=ns$. (c) $s^{n-1}e^{n(1-s)}=\frac{e^{-n\varphi(s)}}s$
and $\mathrm d\varphi=\frac{s-1}s\mathrm ds=v\,\mathrm dv$. For the bounds, compare $\varphi$ with $\frac{(s-1)^2}2$. (d) $n\int_0^\infty v^je^{-nv^2/2}\dd v$ for $j=0,1,2$. (e) The sum is
$\frac{Q(n)+1}{Q(n)+R(n)}$.
\end{hint}

\begin{hint}{pb:dixie}
(a) Given the draws, $P(T_m>N_t)$ averages $P(T_m>k)$ against the Poisson weights; with a Poisson number of draws the $n$ counts are independent Poisson
variables of mean $\frac tn$. (b) Put $y=1-e^{-x}$. (c) Let $x_0=\ln n+(m-1)\ln\ln n$. Above $x_0$ use $1-(1-p)^n\le np$; below, use $1-(1-p)^n\ge1-e^{-np}$.
\end{hint}

\begin{hint}{pb:mappings}
(a) $x,f(x),\dots,f^{k-1}(x)$ distinct and $f^k(x)=x$. (b) Each cycle of length $k$ has $k$ points. (c) Compare $\frac{n^{\underline k}}{n^k}$ with $e^{-k^2/(2n)}$ for $k\le n^{2/3}$, then split
$\sum_k\frac1ke^{-k^2/(2n)}$ into $\sum_{k\le\sqrt n}\frac1k$ and a Riemann sum. You will need $\int_0^\infty(e^{-u}-[u\le1])\frac{\mathrm du}u=-\gamma$.
\end{hint}

\begin{hint}{pb:parking}
(a) On the circle, adding $c$ to every preference rotates the final configuration by $c$. (b) If spot $k+1$ is the first empty one, the cars that preferred a spot at most $k+1$
are exactly the cars parked in $1,\dots,k$. (c) $T$ increases on $\left[0,\frac1e\right[$, and $te^{-t}$ increases from $0$ to $\frac1e$ on $[0,1]$. (d) Stirling.
\end{hint}

\begin{hint}{pb:xplusinv}
Work with $y_n=x_n^2$, which satisfies $y_{n+1}=y_n+2+\frac1{y_n}$. Then $\frac1{y_k}-\frac1{2k+1}=O(\frac{\ln k}{k^2})$, and $\sum_{k=0}^{n-1}\frac1{2k+1}=H_{2n}-\frac12H_n$.
\end{hint}

\begin{hint}{pb:sinesecond}
(a) $\sin^2u=u^2(1-\frac{u^2}3+\frac{2u^4}{45}+O(u^6))$. (b) First sum the crude estimate $\frac1{u_{k+1}^2}-\frac1{u_k^2}=\frac13+O(u_k^2)$, then the precise one. (c) $(\frac n3+a)^{-1/2}$.
\end{hint}

\begin{hint}{pb:tanroots}
(a) $(\tan x-x)'=\tan^2x$. (b) Put $x=a-\varepsilon$; then $\tan\varepsilon=\frac1{a-\varepsilon}$. Get $\varepsilon=O(\frac1a)$ first, then expand $\arctan$.
\end{hint}

\begin{hint}{pb:nbonacci}
(a) $P_{n+1}(x)=xP_n(x)-1$. (b) With $x=2-\varepsilon$, the equation becomes $\varepsilon=(2-\varepsilon)^{-n}$. Get $\varepsilon\le(2-\varepsilon_2)^{-n}$ with $\varepsilon_2=2-r_2$, then iterate.
(c) The characteristic polynomial of the recurrence is $P_n$.
\end{hint}

\begin{hint}{pb:sqrtsum}
In Proposition~11.7 of Tome~4, take $s=-\frac12$ (respectively $s=\frac12$) and $M=2$, and replace $N$ by $n$. Bound the integral remainder.
\end{hint}

\begin{hint}{pb:hyperfactorial}
Differentiate term by term; the factor $(s)_{2k-1}$ vanishes at $s=-1$ for $k\ge2$, so only its derivative survives there. Take $M=3$.
\end{hint}

\begin{hint}{pb:diehit}
(a) Condition on the first roll. (b) On $|z|\le1$, $|D(z)|\le1$, with equality in $D(z)=1$ only if all the $z^j$ are equal to $1$. (c) Subtract the pole at $1$, with residue computed from $D'(1)=\frac72$,
and use Cauchy's estimates on the rest.
\end{hint}

\begin{hint}{pb:banach}
(a) The box found empty is chosen for the $(N+1)$-th time while the other has been chosen $N-r$ times. (b) Sum the relation over $r$. (c) $\binom{2N}N4^{-N}=\frac1{\sqrt{\pi N}}(1+O(\frac1N))$
(Tome~4, Proposition~11.12).
\end{hint}

\begin{hint}{pb:flory}
(a) The pair with the smallest label is filled first and splits the row into two rows, whose labels are still independent and uniform given the first one. (b) Turn the
recurrence into $xV'-V=\frac{2x^2V}{1-x}$, and put $V=xW$. (c) The tail $\sum_{k\ge n}(k-n)\frac{(-2)^k}{k!}$ is alternating.
\end{hint}

\begin{hint}{pb:saw}
(a) Cut a walk of length $m+n$ after $m$ steps. (b) Fix $k$, write $n=qk+r$. (d) After the first step, a self-avoiding walk never goes straight back. For the lower bound, set up a
two-state recurrence (last step horizontal or vertical).
\end{hint}

\begin{hint}{pb:lcs}
(a) Cut both words into a prefix of length $m$ and a suffix of length $n$. (b) Superadditivity gives $\gamma_2\ge\frac{\mathbb E[L_n]}n$ for every $n$. (c) Choose the $k$ positions in each word;
they match with probability $2^{-k}$. Then $\mathbb E[L_n]\le k+nP(L_n\ge k)$.
\end{hint}

\section{Solutions}

\begin{solution}{pb:powersum}
\idea The terms $(1-\frac jn)^n$ are close to $e^{-j}$: the sum is dominated by the last few values of $k$, and the correction comes from the second-order term
of $\ln(1-x)$.

\writeup For $0\le x\le\frac12$, $\ln(1-x)=-\sum_{i\ge1}\frac{x^i}i$, so $\ln(1-x)\le-x-\frac{x^2}2$ and $\ln(1-x)\ge-x-\frac{x^2}2-\frac{x^3}3\sum_{i\ge0}x^i\ge-x-\frac{x^2}2-x^3$. Let
$a_j=(1-\frac jn)^n$, so that the sum is $\sum_{j=0}^{n-1}a_j$. For $j\le\frac n2$, with $x=\frac jn$,
\[
e^{-j}e^{-j^2/(2n)}e^{-j^3/n^2}\le a_j\le e^{-j}e^{-j^2/(2n)}.
\]
Since $1-y\le e^{-y}\le1-y+\frac{y^2}2$ for $y\ge0$, both sides are $e^{-j}\bigl(1-\frac{j^2}{2n}+E_j\bigr)$ with $|E_j|\le\frac{j^4}{8n^2}+\frac{j^3}{n^2}$. So
\[
\sum_{j\le n/2}a_j=\sum_{j\le n/2}e^{-j}-\frac1{2n}\sum_{j\le n/2}j^2e^{-j}+O\Bigl(\frac1{n^2}\sum_jj^4e^{-j}\Bigr).
\]
Replacing the first two sums by the full series costs $O(n^2e^{-n/2})$. For $\frac n2<j<n$, $a_j\le2^{-n}$, so these terms contribute at most $n2^{-n}$. With
$\sum_je^{-j}=\frac e{e-1}$ and $\sum_jj^2x^j=\frac{x(1+x)}{(1-x)^3}$ at $x=\frac1e$, which is $\frac{e(e+1)}{(e-1)^3}$, we get the claimed expansion.

\tools{Tome~0, Chapter~5 (series, $\ln(1-x)$); Tome~4, Chapter~11 (the spirit of Laplace's method).}
\leads Compare with $\sum_kk^k\sim n^n$ (Tome~4, Exercise~11.1): there only the last term matters, here the last $O(1)$ terms all count. The method,
cutting into a region where a Taylor expansion is uniform and a region where everything is exponentially small, is the discrete
Laplace method, and the rest of this chapter uses it again and again.
\end{solution}

\begin{solution}{pb:involasym}
\idea On the circle $|z|=r$ where $|T(z)|r^{-n}$ is smallest along the positive axis, the integrand is a sharp Gaussian bump around $\theta=0$. Laplace's method
does the rest, and Stirling turns the result into a closed form.

\writeup (a) $T$ is entire, so Cauchy's formula on $|z|=r$ (Tome~4, Theorem~9.10) gives $\frac{t_n}{n!}=\frac1{2\pi i}\oint\frac{T(z)}{z^{n+1}}\dd z$; parametrise $z=re^{i\theta}$.

(b) $|T(re^{i\theta})|=\exp\operatorname{Re}\bigl(re^{i\theta}+\frac12r^2e^{2i\theta}\bigr)=\exp\bigl(r\cos\theta+\frac12r^2\cos2\theta\bigr)$, and
$r(\cos\theta-1)+\frac12r^2(\cos2\theta-1)=-2r\sin^2\frac\theta2-r^2\sin^2\theta$.

(c) Let $\delta=r^{-4/5}$. \emph{Tails.} For $\delta\le|\theta|\le\frac\pi2$, $\sin^2\theta\ge\frac{4\theta^2}{\pi^2}$, so by (b) the integrand has modulus at most
$T(r)r^{-n}e^{-4r^2\delta^2/\pi^2}=T(r)r^{-n}e^{-4r^{2/5}/\pi^2}$. For $\frac\pi2\le|\theta|\le\pi$, $\sin^2\frac\theta2\ge\frac12$, so the modulus is at most $T(r)r^{-n}e^{-r}$. Both are
$o\bigl(T(r)r^{-n}/\sqrt b\bigr)$, since $b\le3r^2$.

\emph{Centre.} Let $G(\theta)=re^{i\theta}+\frac12r^2e^{2i\theta}-in\theta$. Then $G(0)=\ln T(r)$, $G'(0)=i(r+r^2)-in=0$, $G''(0)=-(r+2r^2)=-b$, and $|G'''|\le r+4r^2\le5r^2$. By Taylor's formula,
$G(\theta)=\ln T(r)-\frac b2\theta^2+E(\theta)$ with $|E(\theta)|\le r^2|\theta|^3\le r^{-2/5}$ for $|\theta|\le\delta$. Hence
\[
\int_{-\delta}^\delta e^{G(\theta)}\dd\theta=T(r)\bigl(1+O(r^{-2/5})\bigr)\int_{-\delta}^\delta e^{-b\theta^2/2}\dd\theta=T(r)\sqrt{\frac{2\pi}b}\,\bigl(1+o(1)\bigr),
\]
because $b\delta^2\ge2r^{2/5}\to\infty$, so the Gaussian integral over $|\theta|>\delta$ is negligible. Dividing by $2\pi r^n$ gives the claim.

(d) $r=\frac{\sqrt{4n+1}-1}2=\sqrt n-\frac12+\frac1{8\sqrt n}+O(n^{-3/2})$. Since $r^2=n-r$, $\ln T(r)=r+\frac{r^2}2=\frac{n+r}2$. Next $\frac r{\sqrt n}=1+\varepsilon$ with
$\varepsilon=-\frac1{2\sqrt n}+\frac1{8n}+O(n^{-3/2})$, and $\ln(1+\varepsilon)=\varepsilon-\frac{\varepsilon^2}2+O(\varepsilon^3)=-\frac1{2\sqrt n}+O(n^{-3/2})$, so $n\ln r=\frac n2\ln n-\frac{\sqrt n}2+O(n^{-1/2})$. Hence
\[
\ln\frac{T(r)}{r^n}=\frac n2+\frac{\sqrt n}2-\frac14-\frac n2\ln n+\frac{\sqrt n}2+o(1)=\frac n2+\sqrt n-\frac14-\frac n2\ln n+o(1).
\]
With $b\sim2n$ and $n!\sim\sqrt{2\pi n}\,n^ne^{-n}$,
\[
t_n\sim\sqrt{2\pi n}\,n^ne^{-n}\cdot\frac{e^{n/2+\sqrt n-1/4}n^{-n/2}}{\sqrt{4\pi n}}=\frac1{\sqrt2}n^{n/2}e^{-n/2+\sqrt n-1/4}.
\]

\tools{Tome~4, Theorem~9.10 (Cauchy's formula), Theorem~11.9 (Laplace's method) and the example of the central binomial coefficient in Chapter~11; Tome~0, Theorem~5.6 (Stirling).}
\leads This is the saddle-point method of Tome~4 applied to a function that is not a power, the setting that Hayman (1956) axiomatised; the result for
involutions is due to Chowla, Herstein and Moore (1951) and to Moser and Wyman (1955). The notebook shows that the ratio of $t_n$ to the approximation
in (d) is $1.0091$ for $n=1000$: the error is of relative order $\frac1{\sqrt n}$, which is why the approximation in (c), with ratio $0.99983$ at $n=1000$, is much better.
\end{solution}

\begin{solution}{pb:bellasym}
\idea Dobi\'nski's series $\sum\frac{k^n}{k!}$ has a largest term near $k=e^u$, and bounding every term by an exponential in $k$ gives an upper bound that matches the
largest term up to a factor polynomial in $n$.

\writeup (a) The function $x\mapsto x^ne^{-ux}$ on $[0,\infty[$ has derivative $x^{n-1}e^{-ux}(n-ux)$, so it is maximal at $x=\frac nu$, with value $(\frac n{eu})^n$. Hence
$k^n\le(\frac n{eu})^ne^{uk}$ and
\[
b_n\le\frac1e\Bigl(\frac n{eu}\Bigr)^n\sum_k\frac{e^{uk}}{k!}=\Bigl(\frac n{eu}\Bigr)^ne^{e^u-1}.
\]
Since $ue^u=n$, we have $\ln n-\ln u=u$ and $e^u=\frac nu$, so $\ln b_n\le n(u-1)+\frac nu-1$.

(b) Let $m=\lceil e^u\rceil$, so $m=e^u+\eta$ with $0\le\eta<1$, and $b_n\ge\frac{m^n}{e\,m!}$. First $n\ln m\ge n\ln e^u=nu$. Next, Tome~0 (Chapter~5) shows that
$d_m=\ln m!-(m+\frac12)\ln m+m$ decreases from $d_1=1$, so $\ln m!\le m\ln m-m+\frac12\ln m+1$. The function $\psi(x)=x\ln x-x$ has derivative $\ln x$, so
$\psi(m)\le\psi(e^u)+\eta\ln m\le ue^u-e^u+\ln m$. Since $ue^u=n$, $e^u=\frac nu$ and $\ln m\le u+1\le\ln n+1$,
\[
\ln b_n\ge nu-1-\Bigl(n-\frac nu+\frac32\ln m+1\Bigr)\ge n\Bigl(u-1+\frac1u\Bigr)-\frac32\ln n-4.
\]

(c) For $n\ge3$, $u\ge1$, so $u=L-\ln u\le L$, and then $\ln u\le\ln L$, so $u\ge L-\ln L$. Write $\ell=\ln L$. Then $\frac uL=1-\frac{\ln u}L\in\bigl[1-\frac\ell L,1\bigr]$, so
$\ln u=\ell+\ln\frac uL=\ell-\frac{\ln u}L+O(\frac{\ell^2}{L^2})=\ell-\frac\ell L+O(\frac{\ell^2}{L^2})$, and $u=L-\ln u=L-\ell+\frac\ell L+O(\frac{\ell^2}{L^2})$. Also $\frac1u=\frac1L+O(\frac\ell{L^2})$. By (a) and (b),
$\ln b_n=n(u-1+\frac1u)+O(\ln n)$, and
\[
u-1+\frac1u=L-\ell-1+\frac{1+\ell}L+O\Bigl(\frac{\ell^2}{L^2}\Bigr);
\]
the error $O(\ln n)$ is absorbed, since $\frac{\ln n}n=o(\frac{\ell^2}{L^2})$.

\tools{Tome~4, Chapter~2 (Dobi\'nski's formula); Tome~0, Chapter~5 (Stirling's estimate); Tome~0, Chapter~6 (maxima).}
\leads The expansion in (c) converges painfully slowly (for $n=1000$ its main terms give $4400$ against $\ln b_n=4438$), while $n(u-1+\frac1u)$ is off by only $1.92$.
Moser and Wyman (1955) found the next term by the saddle point: $b_n\sim\frac{e^{n(u-1+1/u)-1}}{\sqrt{1+u}}$, and indeed $-1-\frac12\ln(1+u)=-1.916$ for $n=1000$. The
lesson: an implicitly defined quantity like $u$ is often a better scale than its own expansion.
\end{solution}

\begin{solution}{pb:distinctcycles}
\idea The generating function is $\frac1{1-z}$ times a function that is tame at $z=1$; the partial sums of its coefficients then converge to its value at $1$, which
telescopes into $e^{-\gamma}$.

\writeup (a) A permutation is a set of cycles, and there are $(k-1)!$ cycles on a given $k$-set, so the exponential generating function of one cycle of length $k$
is $\frac{z^k}k$. Allowing each length at most once, the exponential formula (Tome~4, Theorem~2.9) gives $\sum_n(\text{number}/n!)z^n=\prod_k(1+\frac{z^k}k)$, and the number divided
by $n!$ is $d_n$. Expanding: $d_1=1$, $d_2=\frac12$ (only $1+\frac{z^2}2$ contributes, not $z\cdot z$), $d_3=\frac13+\frac12=\frac56$, $d_4=\frac14+\frac13=\frac7{12}$, $d_5=\frac15+\frac14+\frac16=\frac{37}{60}$.

(b) For $|z|<1$, $\prod_ke^{z^k/k}=e^{-\ln(1-z)}=\frac1{1-z}$, and all products converge absolutely, so $\prod(1+\frac{z^k}k)=\frac{h(z)}{1-z}$. At $z=1$,
$\prod_{k\le N}(1+\frac1k)e^{-1/k}=(N+1)e^{-H_N}\to e^{-\gamma}$.

(c) $\|fg\|\le\|f\|\,\|g\|$ is the triangle inequality applied to the Cauchy product. Next, $(1+w)e^{-w}=\sum_j\frac{(-w)^j}{j!}+\sum_j\frac{(-1)^jw^{j+1}}{j!}=1-\sum_{j\ge2}(-1)^j\frac{j-1}{j!}w^j$,
so with $w=\frac{z^k}k$,
\[
\Bigl\|\Bigl(1+\frac{z^k}k\Bigr)e^{-z^k/k}-1\Bigr\|\le\sum_{j\ge2}\frac{j-1}{j!}k^{-j}\le\frac1{k^2}\sum_{j\ge2}\frac{j-1}{j!}=\frac1{k^2}.
\]
Let $P_N$ be the product of the first $N$ factors. Then $\|P_N\|\le\prod_k(1+k^{-2})=:C<\infty$, and $\|P_{N+1}-P_N\|\le\|P_N\|\cdot(N+1)^{-2}\le C(N+1)^{-2}$, so $(P_N)$ is a Cauchy
sequence for $\|\cdot\|$ and converges to a series $\tilde h$ with $\|\tilde h\|<\infty$. Convergence in this norm implies convergence of each coefficient and uniform convergence on
$|z|\le1$, so $\tilde h=h$ on $|z|<1$ and $\tilde h(1)=\lim P_N(1)=h(1)$. Hence $\sum_j|h_j|<\infty$ and $\sum_jh_j=h(1)$.

(d) From $\sum d_nz^n=h(z)\sum z^n$, $d_n=h_0+\dots+h_n\to\sum_jh_j=e^{-\gamma}$.

\tools{Tome~4, Theorem~2.9 (exponential formula) and Corollary~11.6 ($H_N-\ln N\to\gamma$); Tome~0, Chapter~5 (absolute convergence, products).}
\leads The convergence is slow: $d_{300}=0.5633$, with an error of order $\frac1n$. Compare with P1, Chapter~4: a permutation with no cycle shorter than $m$ has probability
about $\frac{e^{-\gamma}}m$, and an integer without prime factors below $m$ has density about $\frac{e^{-\gamma}}{\ln m}$ (Mertens). The constant $e^{-\gamma}$ keeps appearing where
cycles or primes are forbidden to repeat or to be small. The norm $\|\cdot\|$ is that of the Wiener algebra, which reappears with Wiener's theorem in Tome~7.
\end{solution}

\begin{solution}{pb:zigzag}
\idea The largest element splits an alternating permutation into two alternating pieces; this gives a Riccati equation, solved by $\sec+\tan$. The pole at $\frac\pi2$ then
dictates the growth, and the next pole at $-\frac{3\pi}2$ the error.

\writeup (a) Call a permutation of $m\ge2$ elements a zigzag if it is alternating starting up ($w_1<w_2>\cdots$) or starting down ($w_1>w_2<\cdots$). Complementing
($w_i\mapsto m+1-w_i$) exchanges the two kinds, so there are $2A_m$ zigzags. Let $w$ be a zigzag of $\{1,\dots,n+1\}$, $n\ge1$, and let $n+1$ be at position $k+1$. It is a peak, so
the block to its left, read from right to left, starts with a descent from $n+1$ and then goes up, down, \dots: reversed, it is an alternating permutation
(starting up) of a $k$-set, or has length at most $1$. Likewise the block to the right is alternating starting up. Conversely any choice of the $k$-set, of an alternating
arrangement of it (reversed on the left) and of an alternating arrangement of the rest gives a zigzag, of a kind determined by the parity of $k$. So
$\sum_k\binom nkA_kA_{n-k}=2A_{n+1}$.

(b) The coefficient of $\frac{z^n}{n!}$ in $f^2$ is $\sum_k\binom nkA_kA_{n-k}$, in $2f'$ it is $2A_{n+1}$; by (a) they agree for $n\ge1$, and for $n=0$, $2A_1=2=1+A_0^2$. So $2f'=1+f^2$
as formal series. The function $g=\frac{1+\sin z}{\cos z}=\sec z+\tan z$ satisfies $g(0)=1$ and
\[
2g'=2\,\frac{\cos^2z+\sin z(1+\sin z)}{\cos^2z}=\frac{2(1+\sin z)}{\cos^2z}=\frac{(1+\sin z)^2+\cos^2z}{\cos^2z}=1+g^2.
\]
The equation $2f'=1+f^2$ with $f(0)=1$ determines the Taylor coefficients one after another, so $f$ is the Taylor series of $g$, which converges on $|z|<\frac\pi2$.

(c) The zeros of $\cos$ are $\frac\pi2+j\pi$, all simple. At $z=\frac\pi2$, $1+\sin z=2$ and $(\cos)'=-1$, so $g$ has a simple pole with residue $-2$, that is principal part
$\frac{-2}{z-\pi/2}=\frac4{\pi-2z}$. At $z=-\frac\pi2$, $1+\sin z$ has a double zero, so the singularity is removable. So $g-\frac4{\pi-2z}$ is holomorphic on $|z|<\frac{3\pi}2$. By Cauchy's
estimates (Tome~4, Corollary~9.12) on a circle of radius $\rho<\frac{3\pi}2$, its coefficients are $O(\rho^{-n})$. Since $[z^n]\frac4{\pi-2z}=\frac4\pi(\frac2\pi)^n=2(\frac2\pi)^{n+1}$, the relative error is
$O\bigl((\frac\pi{2\rho})^n\bigr)$, and $\frac\pi{2\rho}$ can be any number greater than $\frac13$.

(d) $\frac{A_{10}}{10!}=\frac{50521}{3628800}=0.0139222$, and $2(\frac2\pi)^{11}=0.0139223$: the relative error is $-5.6\cdot10^{-6}$, close to $-3^{-11}$.

\tools{Tome~4, Chapter~2 (exponential generating functions), Chapter~9 (poles, Cauchy's estimates).}
\leads D\'esir\'e Andr\'e found (b) in 1879. Subtracting all the poles gives the exact formula
$\frac{A_n}{n!}=2(\frac2\pi)^{n+1}\sum_{k\ge0}(-1)^{k(n+1)}(2k+1)^{-(n+1)}$, so the error is about $(-1)^{n+1}3^{-(n+1)}$, as observed in (d). For odd $n$ the sum is
$(1-2^{-n-1})\zeta(n+1)$: the tangent numbers $A_{2m-1}$ are the coefficients behind $\zeta(2m)$, which Tome~4, Chapter~10 computes with Bernoulli numbers.
\end{solution}

\begin{solution}{pb:binompowers}
\idea Raising the binomial law to the power $s$ keeps it Gaussian, with variance divided by $s$. Near the middle the local limit theorem applies; away from it the terms
are too small to matter.

\writeup Let $p_k=\binom nk2^{-n}$, $x_k=\frac{2k-n}{\sqrt n}$ and $\phi(x)=\frac{e^{-x^2/2}}{\sqrt{2\pi}}$, so the sum is $2^{sn}\Sigma$ with $\Sigma=\sum_kp_k^s$. Fix $A>0$. By Tome~3, Theorem~6.12 (with $p=\frac12$,
$\sqrt{npq}=\frac{\sqrt n}2$), $p_k=\frac2{\sqrt n}\phi(x_k)(1+\varepsilon_k)$ with $\max_{|x_k|\le A}|\varepsilon_k|\to0$. The points $x_k$ are spaced $\frac2{\sqrt n}$ apart, so by Riemann sums
\[
\Bigl(\frac{\sqrt n}2\Bigr)^{s-1}\sum_{|x_k|\le A}p_k^s=\bigl(1+o(1)\bigr)\frac2{\sqrt n}\sum_{|x_k|\le A}\phi(x_k)^s\to\int_{-A}^A\phi^s.
\]
For the rest, $p_k\le p_{\lfloor n/2\rfloor}\le\frac c{\sqrt n}$ (central binomial coefficient), and $\sum_{|x_k|>A}p_k\le\frac1{A^2}$ by Chebyshev, since $x_K$ has mean $0$ and variance $1$ when $K$ is binomial.
So $(\frac{\sqrt n}2)^{s-1}\sum_{|x_k|>A}p_k^s\le(\frac c2)^{s-1}A^{-2}$. Letting $n\to\infty$, then $A\to\infty$,
\[
\Bigl(\frac{\sqrt n}2\Bigr)^{s-1}\Sigma\to\int_\R\phi^s=(2\pi)^{-s/2}\sqrt{\frac{2\pi}s}=\frac{(2\pi)^{(1-s)/2}}{\sqrt s}.
\]
Hence $\Sigma\sim(\frac2{\sqrt n})^{s-1}\frac{(2\pi)^{(1-s)/2}}{\sqrt s}=\frac1{\sqrt s}(\frac2{\pi n})^{(s-1)/2}$. For $s=2$ this gives $\frac{4^n}{\sqrt2}\sqrt{\frac2{\pi n}}=\frac{4^n}{\sqrt{\pi n}}$, which is $\binom{2n}n$ asymptotically (by
Vandermonde, $\sum\binom nk^2=\binom{2n}n$). For $s=3$: $\sum\binom nk^3\sim\frac{2\cdot8^n}{\sqrt3\,\pi n}$; for $s=4$: $\sum\binom nk^4\sim\frac{16^n}2\bigl(\frac2{\pi n}\bigr)^{3/2}$. The notebook finds ratios $0.9997$
and $0.9994$ at $n=1000$.

\tools{Tome~3, Theorem~6.12 (local de Moivre--Laplace) and Chapter~4 (Chebyshev); Tome~4, Chapter~4 (Gaussian integral).}
\leads For $s=3$ these are the Franel numbers; Franel (1894, 1895) showed that the sums for $s=3$ and $s=4$ satisfy linear recurrences with polynomial coefficients, and this
asymptotic is what such a recurrence must reproduce. The same local-limit argument gives $\sum_kP(S_n=k)^s$ for any lattice random walk.
\end{solution}

\begin{solution}{pb:poissonise}
\idea Words with bounded letter counts are counted by $e_m(x)^n$, and the factor $k!$ is an integral: summing over $k$ turns the generating function into a Laplace
integral, which then has an easy asymptotic.

\writeup (a) $T>k$ means that in the first $k$ draws every value appears fewer than $m$ times. The words of length $k$ on $n$ letters with each letter used at most $m-1$
times number $k![x^k]e_m(x)^n$ (choose the multiplicities $j_1+\dots+j_n=k$, then $\frac{k!}{j_1!\cdots j_n!}$ arrangements), and each word has probability $n^{-k}$.

(b) Write $e_m(x)^n=\sum_ka_kx^k$, a polynomial. Then $\mathbb E[T]=\sum_{k\ge0}P(T>k)=\sum_ka_kk!\,n^{-k}=\sum_ka_kn^{-k}\int_0^\infty t^ke^{-t}\dd t=\int_0^\infty e^{-t}e_m(t/n)^n\dd t$.

(c) For $m=2$, $\int_0^\infty e^{-t}(1+\frac tn)^n\dd t=\sum_j\binom nj\frac{j!}{n^j}=\sum_{j=0}^n\frac{n^{\underline j}}{n^j}=1+Q(n)$. Tome~3, Exercise~2.7 is the same sum, $\sum_kq_k$.

(d) Put $t=nx$: $\mathbb E[T]=n\int_0^\infty\psi(x)^n\dd x$ with $\psi(x)=e^{-x}(1+x+\frac{x^2}2)$. Then $\psi(0)=1$ and $\psi'(x)=-\frac{x^2}2e^{-x}<0$, so $\psi$ decreases, and
$\ln\psi(x)=-\frac{x^3}6+O(x^4)$. Choose $\delta>0$ with $-\frac{x^3}6-Cx^4\le\ln\psi(x)\le-\frac{x^3}{12}$ on $[0,\delta]$. On $[\delta,\infty[$, $\psi^n\le\psi(\delta)^{n-1}\psi$ and $\int\psi=3$, which is exponentially
small. On $[0,\delta]$ put $x=(6/n)^{1/3}y$: then $n\ln\psi(x)\to-y^3$ with domination by $e^{-y^3/2}$, so
\[
n\int_0^\delta\psi^n\dd x=n\Bigl(\frac6n\Bigr)^{1/3}\int_0^\infty e^{-y^3}\dd y\,\bigl(1+o(1)\bigr)=(6n^2)^{1/3}\,\Gamma\Bigl(\frac43\Bigr)\bigl(1+o(1)\bigr),
\]
using $\int_0^\infty e^{-y^3}\dd y=\Gamma(\frac43)$ (Tome~4, Exercise~5.2).

(e) $\Gamma(\frac43)(6\cdot365^2)^{1/3}=82.9$, against the exact $88.74$. So the answer to the think box is ``about $90$''. (For pairs, $\sqrt{\pi\cdot365/2}=23.9$ against $24.62$.)

\tools{Tome~3, Exercise~2.7 (waiting for a pair); Tome~4, Chapter~2 (exponential generating functions), Chapter~5 (Gamma) and Chapter~11 (Laplace's method at an end point).}
\leads Klamkin and Newman (1967) found the general answer $\mathbb E[T]\sim(m!)^{1/m}\Gamma(1+\frac1m)\,n^{1-1/m}$ (for $m=2$ it is $\sqrt{\pi n/2}$). The relative error
decays only like $n^{-1/3}$ for $m=3$: $1\%$ at $n=10^5$. Hash tables use the same computation, with $m$ the capacity of a bucket.
\end{solution}

\begin{solution}{pb:ramanujan}
\idea Both $Q$ and $R$ are Laplace integrals of the same function, on the two sides of its peak. The peak is not symmetric: the skewness of $s^{n-1}e^{n(1-s)}$, seen through
$g(v)=1-\frac v3+\cdots$, gives $\mp\frac13$, and Ramanujan's $\frac13$ is half of $\frac13+\frac13$.

\writeup (a) $\frac{n^{n-j}}{(n-j)!}=\frac{n^n}{n!}\cdot\frac{n(n-1)\cdots(n-j+1)}{n^j}$, and summing over $j=1,\dots,n$ gives the first identity. Likewise $\frac{n^{n+j}}{(n+j)!}=\frac{n^n}{n!}\cdot\frac{n^j}{(n+1)\cdots(n+j)}$, summed
over $j\ge0$. Adding gives $e^n$. Then $\theta_n\frac{n^n}{n!}=\frac{e^n}2-\frac{n^n}{n!}Q=\frac{n^n}{n!}\bigl(\frac{Q+R}2-Q\bigr)$.

(b) The function $F(x)=\sum_{k\ge n}e^{-x}\frac{x^k}{k!}$ satisfies $F(0)=0$ and $F'(x)=\sum_{k\ge n}e^{-x}\bigl(\frac{x^{k-1}}{(k-1)!}-\frac{x^k}{k!}\bigr)=e^{-x}\frac{x^{n-1}}{(n-1)!}$ (telescoping). So
$\sum_{k\ge n}\frac{n^k}{k!}=\frac{e^n}{(n-1)!}\int_0^nt^{n-1}e^{-t}\dd t$, and with $t=ns$ and (a),
\[
R(n)=\frac{n!}{n^n}\cdot\frac{e^n}{(n-1)!}\,n^n\int_0^1s^{n-1}e^{-ns}\dd s=n\int_0^1s^{n-1}e^{n(1-s)}\dd s.
\]
The same computation with $\int_0^\infty t^{n-1}e^{-t}\dd t=(n-1)!$ gives $Q+R=n\int_0^\infty s^{n-1}e^{n(1-s)}\dd s$, hence the formula for $Q$.

(c) $\varphi$ is convex with minimum $\varphi(1)=0$, so $v$ is an increasing bijection, smooth with $v'(1)=1$ (since $\varphi(s)=\frac{(s-1)^2}2+O((s-1)^3)$), and its inverse is smooth. Since
$s^{n-1}e^{n(1-s)}=\frac{e^{-n\varphi(s)}}s$ and $v\,\mathrm dv=\mathrm d\varphi=\frac{s-1}s\mathrm ds$, we get $\frac{\mathrm ds}s=g(v)\,\mathrm dv$, which gives the two integrals. $g>0$ since $v$ and $s-1$ have the same sign.
For $s>1$, $\ln s\ge(s-1)-\frac{(s-1)^2}2$, so $\varphi\le\frac{(s-1)^2}2$, $v\le s-1$ and $g\le1$. For $s<1$ we need $|v|\le(1-s)(1+|v|)$, that is $|v|\le\frac{1-s}s=:w$; with $s=\frac1{1+w}$,
$2\varphi=2\ln(1+w)-\frac{2w}{1+w}$, and $k(w)=w^2-2\ln(1+w)+\frac{2w}{1+w}$ has $k(0)=0$ and $k'(w)=2w-\frac{2w}{(1+w)^2}\ge0$, so $2\varphi\le w^2$. Next, with $x=s-1$, $\varphi=\frac{x^2}2-\frac{x^3}3+O(x^4)$,
so $v=x(1-\frac x3+O(x^2))$, $x=v+\frac{v^2}3+O(v^3)$ and $g=\frac vx=1-\frac v3+O(v^2)$ near $0$. As $g$ is smooth, $|g(v)-1+\frac v3|\le Cv^2$ on $[-1,1]$; for $|v|\ge1$,
$|g(v)-1+\frac v3|\le2+2|v|\le4v^2$.

(d) With $I_j=n\int_0^\infty v^je^{-nv^2/2}\dd v$: $I_0=\sqrt{\frac{\pi n}2}$, $I_1=1$, $I_2=\frac{\sqrt{2\pi}}{2\sqrt n}$. So $Q=I_0-\frac13I_1+O(I_2)=\sqrt{\frac{\pi n}2}-\frac13+O(n^{-1/2})$, and the same on
$\left]-\infty,0\right]$ gives $R=\sqrt{\frac{\pi n}2}+\frac13+O(n^{-1/2})$. Hence $\theta_n=\frac{R-Q}2=\frac13+O(n^{-1/2})$.

(e) $e^{-n}\sum_{k\le n}\frac{n^k}{k!}=e^{-n}\frac{n^n}{n!}(Q+1)=\frac{Q+1}{Q+R}$. Now $Q+R=\frac{n!e^n}{n^n}=\sqrt{2\pi n}\,(1+O(\frac1n))$ (Tome~4, Theorem~11.8) and $Q+1=\frac12\sqrt{2\pi n}+\frac23+O(n^{-1/2})$, so the ratio is
$\frac12+\frac2{3\sqrt{2\pi n}}+O(\frac1n)$.

\tools{Tome~3, Exercise~6.5 (the limit $\frac12$); Tome~4, Theorem~11.8 (Stirling's series) and Theorem~11.9 (Laplace's method); Tome~0, Chapter~6 (inverse functions).}
\leads The notebook finds $n(\theta_n-\frac13)\to0.029630$, which is $\frac4{135}$: the next term of the expansion, which comes from the $v^2$ term of $g$. The same Laplace
integral, cut at other places, gives Knuth's analysis of linear probing: with a table of $m$ cells and $n$ keys, the cost of a search involves $Q$-like
sums at the ratio $\frac nm$, and the full table is exactly Ramanujan's case.
\end{solution}

\begin{solution}{pb:dixie}
\idea With a Poisson number of draws the coupon counts are independent, so ``every type seen $m$ times'' has a product probability; the integral is then a matter of
locating where $np_m(x)$ crosses $1$.

\writeup (a) For each $t$, $P(T_m>N_t)=\sum_kP(T_m>k)e^{-t}\frac{t^k}{k!}$, and $\int_0^\infty e^{-t}\frac{t^k}{k!}\dd t=1$; by Tonelli, $\int_0^\infty P(T_m>N_t)\dd t=\sum_kP(T_m>k)=\mathbb E[T_m]$. Given $N_t$,
the counts of the $n$ types are multinomial; unconditionally they are independent Poisson variables of mean $\frac tn$, since with $c=c_1+\dots+c_n$,
$P(\text{counts}=(c_i))=e^{-t}\frac{t^c}{c!}\cdot\frac{c!}{c_1!\cdots c_n!}n^{-c}=\prod_ie^{-t/n}\frac{(t/n)^{c_i}}{c_i!}$. So
$P(T_m\le N_t)=(1-p_m(t/n))^n$, and $t=nx$ gives the formula.

(b) For $m=1$, $p_1(x)=e^{-x}$; with $y=1-e^{-x}$, $\mathrm dx=\frac{\mathrm dy}{1-y}$ and $\int_0^\infty(1-(1-e^{-x})^n)\dd x=\int_0^1\frac{1-y^n}{1-y}\dd y=\sum_{j<n}\frac1{j+1}=H_n$.

(c) Write $p=p_m$, $F(x)=1-(1-p(x))^n$, $x_0=\ln n+(m-1)\ln\ln n$; we show $\int_0^\infty F=x_0+O(1)$. \emph{Upper bound.} $\int_0^{x_0}F\le x_0$. For $x\ge x_0$ (assume $x_0\ge1$),
$F\le np(x)\le nmx^{m-1}e^{-x}$, since each term $\frac{x^j}{j!}$ with $j<m$ is at most $x^{m-1}$. With $x=x_0+y$, $x^{m-1}\le x_0^{m-1}(1+y)^{m-1}$ and
$nx_0^{m-1}e^{-x_0}=(\frac{x_0}{\ln n})^{m-1}\le2^{m-1}$ for large $n$, so $\int_{x_0}^\infty F\le m2^{m-1}\int_0^\infty(1+y)^{m-1}e^{-y}\dd y=O(1)$.
\emph{Lower bound.} $1-F=(1-p)^n\le e^{-np}$. For $x=x_0-y\in[\frac{x_0}2,x_0]$, $np(x)\ge ne^{-x}\frac{x^{m-1}}{(m-1)!}\ge e^y\cdot ne^{-x_0}\frac{(x_0/2)^{m-1}}{(m-1)!}\ge ce^y$ with $c>0$ independent of $n$, so
$\int_{x_0/2}^{x_0}(1-F)\le\int_0^\infty e^{-ce^y}\dd y=O(1)$. For $x\le\frac{x_0}2$, $np(x)\ge ne^{-x}\ge ne^{-x_0/2}=\sqrt n(\ln n)^{-(m-1)/2}$, so $\int_0^{x_0/2}(1-F)\le x_0e^{-\sqrt n(\ln n)^{-(m-1)/2}}\to0$.
Hence $\int_0^\infty F\ge x_0-O(1)$, and $\mathbb E[T_m]=nx_0+O(n)$.

\tools{Tome~3, Theorem~2.14 (coupon collector) and Chapter~2 (Poisson law); Tome~0, Chapter~6 (integrals).}
\leads Newman and Shepp (1960) found these two terms; Erd\H os and R\'enyi (1961) found the constant: $\mathbb E[T_m]=n\ln n+(m-1)n\ln\ln n+n(\gamma-\ln(m-1)!)+o(n)$,
together with a Gumbel limit law (for $m=1$ see Tome~3, Chapter~9). The $o(n)$ is deceptive: for $m=2$ and $n=10^6$ the notebook still finds a gap of $0.28\,n$.
\end{solution}

\begin{solution}{pb:mappings}
\idea Following $x,f(x),f^2(x),\dots$ is drawing with replacement until a repeat, so the birthday sums reappear; the number of cycles weights the same terms
by $\frac1k$, which turns $\sqrt n$ into a logarithm.

\writeup (a) $x$ lies on a cycle of length exactly $k$ if and only if $x,f(x),\dots,f^{k-1}(x)$ are distinct and $f^k(x)=x$. The values $f(x),f(f(x)),\dots$ are fresh uniform draws as long as
no repeat occurs, so this has probability $\frac{n-1}n\cdots\frac{n-k+1}n\cdot\frac1n=\frac{n^{\underline k}}{n^{k+1}}$. Summing over $k$ and over the $n$ points gives $\mathbb E[\#\text{cyclic}]=\sum_k\frac{n^{\underline k}}{n^k}=Q(n)$,
and $Q(n)\sim\sqrt{\pi n/2}$ by Problem~\ref{pb:ramanujan} (or Tome~3, Exercise~2.7).

(b) The number of cycles is $\sum_x\frac{[x\text{ cyclic}]}{\text{length of its cycle}}$, so $C_n=n\sum_k\frac1k\cdot\frac{n^{\underline k}}{n^{k+1}}$.

(c) Let $w_k=\frac{n^{\underline k}}{n^k}$. \emph{Comparison.} $\ln w_k=\sum_{j<k}\ln(1-\frac jn)$. For $k\le n^{2/3}$ and $n\ge8$, $\frac jn\le\frac12$, and $-x-x^2\le\ln(1-x)\le-x$ gives
$\bigl|\ln w_k+\frac{k^2}{2n}\bigr|\le\frac k{2n}+\frac{k^3}{n^2}\le2$, hence $|w_k-e^{-k^2/(2n)}|\le Ce^{-k^2/(2n)}(\frac kn+\frac{k^3}{n^2})$. Then
$\sum_k\frac1ke^{-k^2/(2n)}(\frac kn+\frac{k^3}{n^2})=O(\frac{\sqrt n}n)+O(\frac{n^{3/2}}{n^2})\to0$. For $k>n^{2/3}$, $w_k\le w_{\lceil n^{2/3}\rceil}\le e^{-n^{1/3}/3}$ and $e^{-k^2/(2n)}\le e^{-n^{1/3}/2}$, so these terms total
at most $2ne^{-n^{1/3}/3}\to0$. So $C_n=S_n+o(1)$ with $S_n=\sum_{k\ge1}\frac1ke^{-k^2/(2n)}$.

\emph{The sum.} With $a=n^{-1/2}$ and $\chi(x)=\frac{e^{-x^2/2}-[x\le1]}x$, $S_n=\sum_{k\le\sqrt n}\frac1k+\sum_{k\ge1}a\chi(ka)$. The first sum is $\frac12\ln n+\gamma+o(1)$ (Tome~4, Corollary~11.6). The function
$\chi$ is bounded on $\left]0,1\right]$, monotone on $\left]0,1\right]$ and on $\left]1,\infty\right[$, and integrable, so its Riemann sums converge: $\sum_ka\chi(ka)\to\int_0^\infty\chi$. With $u=\frac{x^2}2$,
\[
\int_0^\infty\chi=\frac12\int_0^\infty\bigl(e^{-u}-[u\le\tfrac12]\bigr)\frac{\mathrm du}u=\frac12\Bigl(\int_0^\infty\bigl(e^{-u}-[u\le1]\bigr)\frac{\mathrm du}u+\int_{1/2}^1\frac{\mathrm du}u\Bigr)=\frac12(-\gamma+\ln2).
\]
The value $-\gamma$ comes from integrating by parts with $\ln u$, which turns the integral into $\int_0^\infty e^{-u}\ln u\dd u=\Gamma'(1)=-\gamma$ (Tome~4, Exercise~5.7). So
$C_n=\frac12\ln n+\gamma+\frac12(\ln2-\gamma)+o(1)=\frac12(\ln2n+\gamma)+o(1)$.

\tools{Tome~3, Chapter~1 (birthdays); Tome~4, Corollary~11.6 (harmonic numbers) and Exercise~5.7 ($\Gamma'(1)=-\gamma$).}
\leads Kruskal (1954) found (c), and Harris (1960) the distributions. The rho shape of the path $x,f(x),f^2(x),\dots$ (a tail then a cycle, both of length about
$\sqrt n$) is what Pollard's factoring method exploits. The notebook finds $C_n-\frac12(\ln2n+\gamma)\approx\frac{0.42}{\sqrt n}$.
\end{solution}

\begin{solution}{pb:parking}
\idea On a circle nobody fails to park and symmetry is perfect; cutting the circle at the empty spot gives the street. On the street, the first empty spot splits the
cars into a parking function and an unconstrained rest.

\writeup (a) On the circle of $n+1$ spots every car parks, since it drives around until it finds a free spot, and exactly one spot stays empty. Adding $c$ (mod $n+1$) to every
preference moves every car's spot by $c$, so it moves the empty spot by $c$. The $(n+1)^n$ preference sequences therefore fall into classes of $n+1$ rotations,
in each of which exactly one sequence leaves spot $n+1$ empty. A sequence leaves $n+1$ empty if and only if no car ever reaches spot $n+1$, that is, all its preferences are in
$\{1,\dots,n\}$ and every car parks in $1,\dots,n$ on the street: it is a parking function. So there are $\frac{(n+1)^n}{n+1}=(n+1)^{n-1}$.

(b) Suppose spot $k+1$ is the first empty one at the end. A car preferring a spot $\le k+1$ parks at the first free spot from its preference, which is at most $k+1$ since
$k+1$ stays free; it is not $k+1$, so it is in $\{1,\dots,k\}$. A car preferring a spot $\ge k+2$ parks beyond $k+1$. So the cars parked in $1,\dots,k$ are exactly the
$k$ cars with preference $\le k+1$; none prefers $k+1$ (it would park there), and their preferences form a parking function of length $k$. The other $n-k$ cars have
preferences in $\{k+2,\dots,n+1\}$, freely. Conversely such choices produce first empty spot $k+1$. This gives $\binom nk(k+1)^{k-1}(n-k)^{n-k}$ sequences, and summing over $k$ gives
$(n+1)^n$.

(c) $\binom nk\frac{(n-k)^{n-k}}{(n+1)^n}=\frac{n^{\underline k}}{n^kk!}\cdot\frac{(1-k/n)^{n-k}}{(1+1/n)^n}\to\frac{e^{-k}}{k!e}$. With $m=k+1$, $p_k=\frac{m^{m-2}e^{-m}}{(m-1)!}=\frac{m^{m-1}}{m!}e^{-m}$, so $\sum_kp_k=\sum_{m\ge1}\frac{m^{m-1}}{m!}e^{-m}$. The series
$T(x)$ has radius $\frac1e$ (Stirling), is continuous and increasing on $\left[0,\frac1e\right[$ with $T(0)=0$, and $T(x)e^{-T(x)}=x$. The map $t\mapsto te^{-t}$ increases from $0$ to $\frac1e$ on $[0,1]$, and
$T(x)$ never equals $1$ for $x<\frac1e$, so $T(x)<1$ and $T(x)$ is the inverse of $te^{-t}$ there; as $x\to\frac1e$, $T(x)\to1$. By monotone convergence, $\sum_m\frac{m^{m-1}}{m!}e^{-m}=1$.

(d) By Stirling, $\frac{m^{m-1}}{m!}e^{-m}\sim\frac{m^{m-1}e^{-m}}{\sqrt{2\pi m}\,m^me^{-m}}=\frac1{\sqrt{2\pi}m^{3/2}}$, and $m\sim k$. So $\sum_kkp_k=\infty$.

\tools{the tree-function identity $T=ze^T$; Tome~0, Theorem~5.6 (Stirling).}
\leads Konheim and Weiss (1966) introduced parking functions to analyse linear probing; the circular argument is Pollak's, recorded by Riordan (1969). The limit law
$(p_k)$ is the Borel distribution: $K+1$ has the law of the total size of a Galton--Watson tree with Poisson offspring of mean $1$, a critical tree, hence the heavy
tail. For finite $n$, the notebook finds $\mathbb E[K]\approx1.23\sqrt n$ at $n=3000$.
\end{solution}

\begin{solution}{pb:xplusinv}
\idea Squaring linearises the recursion: $x_n^2$ grows by $2$ plus a small term $\frac1{x_n^2}$, which is about $\frac1{2n}$ and so adds up to a logarithm. Feeding the crude estimate back
gives the constant.

\writeup Let $y_n=x_n^2$; then $y_{n+1}=y_n+2+\frac1{y_n}$ and $y_0=1$.

(a) $y_{n+1}\ge y_n+2$ gives $y_n\ge2n+1$. Then $y_n=2n+1+\sum_{k=0}^{n-1}\frac1{y_k}\le2n+1+\sum_{k=0}^{n-1}\frac1{2k+1}\le2n+2+\int_0^{n-1}\frac{\mathrm dx}{2x+1}=2n+2+\frac12\ln(2n-1)$, which is at most
$2n+2+\ln n$ for $n\ge1$.

(b) Let $\delta_k=\frac1{y_k}-\frac1{2k+1}=\frac{2k+1-y_k}{y_k(2k+1)}$. By (a), $-\frac{1+\ln k}{(2k+1)^2}\le\delta_k\le0$ for $k\ge1$, and $\delta_0=0$, so $\sum_k\delta_k$ converges absolutely to some $\Delta$, and
$\sum_{k<n}\delta_k=\Delta+o(1)$. Also $\sum_{k=0}^{n-1}\frac1{2k+1}=H_{2n}-\frac12H_n=\ln2n+\gamma-\frac12(\ln n+\gamma)+o(1)=\frac12\ln n+\ln2+\frac\gamma2+o(1)$ (Tome~4, Corollary~11.6). So
\[
y_n=2n+\frac12\ln n+C+o(1),\qq C=1+\ln2+\frac\gamma2+\Delta.
\]

(c) With $a_n=\frac12\ln n+C+o(1)=O(\ln n)$, $\sqrt{2n+a_n}=\sqrt{2n}\bigl(1+\frac{a_n}{4n}+O(\frac{a_n^2}{n^2})\bigr)=\sqrt{2n}+\frac{a_n}{2\sqrt{2n}}+O\bigl(\frac{\ln^2n}{n^{3/2}}\bigr)$, and $\frac{a_n}{2\sqrt{2n}}=\frac{\ln n+2C}{4\sqrt{2n}}+o(\frac1{\sqrt n})$.

\tools{Tome~0, Chapter~5 (sequences, comparison with integrals); Tome~4, Corollary~11.6 (harmonic numbers).}
\leads Here $\Delta=-0.2586\ldots$, so $C=1.72314\ldots$; the constant depends on $x_0$, and we know no closed form for it. The same squaring trick handles $x_{n+1}=x_n+\frac1{x_n^p}$ (raise to the
power $p+1$). De Bruijn's \emph{Asymptotic Methods in Analysis} (1958) develops this bootstrapping for iterations in general.
\end{solution}

\begin{solution}{pb:sinesecond}
\idea The quantity $\frac1{u_n^2}$ grows by almost exactly $\frac13$ at each step; the defect $\frac{u_n^2}{15}$ is about $\frac1{5n}$, so it adds up to $\frac15\ln n$. Inverting gives the second term
of $u_n$.

\writeup (a) $\sin u=u\bigl(1-\frac{u^2}6+\frac{u^4}{120}+O(u^6)\bigr)$, so $\sin^2u=u^2\bigl(1-\frac{u^2}3+\bigl(\frac1{36}+\frac1{60}\bigr)u^4+O(u^6)\bigr)=u^2\bigl(1-\varepsilon\bigr)$ with $\varepsilon=\frac{u^2}3-\frac{2u^4}{45}+O(u^6)$. Then
$\frac1{1-\varepsilon}=1+\varepsilon+\varepsilon^2+O(\varepsilon^3)=1+\frac{u^2}3+\bigl(\frac19-\frac2{45}\bigr)u^4+O(u^6)=1+\frac{u^2}3+\frac{u^4}{15}+O(u^6)$; divide by $u^2$.

(b) Let $e_k=\frac1{u_{k+1}^2}-\frac1{u_k^2}-\frac13$. The function $\frac1{\sin^2u}-\frac1{u^2}-\frac13$ is continuous on $\left]0,\frac\pi2\right]$ and $O(u^2)$ at $0$, so $|e_k|\le c\,u_k^2$. By Tome~0,
Exercise~5.4, $\frac1{u_{k+1}^2}-\frac1{u_k^2}\to\frac13$, so $\frac1{u_k^2}\ge\frac k4$ for $k\ge k_0$, and there is $c_1$ with $u_k^2\le\frac{c_1}{k+1}$ for all $k$. Hence $|\sum_{k<n}e_k|\le cc_1\sum_{k<n}\frac1{k+1}=O(\ln n)$ and
$\frac1{u_n^2}=\frac1{u_0^2}+\frac n3+\sum_{k<n}e_k=\frac n3+O(\ln n)$. Therefore $u_k^2=\frac3k+O(\frac{\ln k}{k^2})$, and by (a) $e_k=\frac{u_k^2}{15}+O(u_k^4)=\frac1{5k}+O(\frac{\ln k}{k^2})$. Summing,
$\sum_{k=1}^{n-1}e_k=\frac15\ln n+c'+o(1)$, and $\frac1{u_n^2}=\frac n3+\frac15\ln n+C+o(1)$.

(c) With $a=\frac15\ln n+C+o(1)$, $u_n=\bigl(\frac n3+a\bigr)^{-1/2}=\sqrt{\frac3n}\bigl(1+\frac{3a}n\bigr)^{-1/2}=\sqrt{\frac3n}\bigl(1-\frac{3a}{2n}+O(\frac{a^2}{n^2})\bigr)$, and $\sqrt{\frac3n}\cdot\frac3{2n}\cdot\frac{\ln n}5=\frac{3\sqrt3}{10}\cdot\frac{\ln n}{n^{3/2}}$.
The term in $C$ is $O(n^{-3/2})$.

\tools{Tome~0, Exercise~5.4 (the first term) and Chapter~6 (Taylor expansions).}
\leads The correction is visible but slow: the notebook finds $\frac1{u_n^2}-\frac n3$ equal to $2.19$ at $n=10^3$ and $3.57$ at $n=10^6$, increasing by $\frac15\ln10=0.46$ for each factor $10$. For any
$u_{n+1}=u_n-au_n^p+bu_n^{2p-1}+\cdots$ the same computation gives a main term in $n^{-1/(p-1)}$ and a logarithmic correction, unless a coefficient happens to vanish.
\end{solution}

\begin{solution}{pb:tanroots}
\idea Near $a=(k+\frac12)\pi$ the tangent is huge, so $\tan x=x$ forces $x$ just below $a$; writing $x=a-\varepsilon$ turns the equation into $\tan\varepsilon=\frac1{a-\varepsilon}$, which can be iterated.

\writeup (a) On $I_k=\left]k\pi-\frac\pi2,k\pi+\frac\pi2\right[$ the function $h(x)=\tan x-x$ has $h'=\tan^2x\ge0$, vanishing only at $k\pi$, so $h$ is increasing, from $-\infty$ to $+\infty$: exactly one root.
Since $h(k\pi)=-k\pi<0$, the root is in $\left]k\pi,k\pi+\frac\pi2\right[$.

(b) Put $x_k=a-\varepsilon$ with $0<\varepsilon<\frac\pi2$. Then $\tan(a-\varepsilon)=\cot\varepsilon$, so the equation is $\tan\varepsilon=\frac1{a-\varepsilon}$. Since $a-\varepsilon>a-\frac\pi2\ge\frac a3$, $\varepsilon<\tan\varepsilon\le\frac3a$. Then
$\frac1{a-\varepsilon}=\frac1a\bigl(1+\frac\varepsilon a+O(\frac{\varepsilon^2}{a^2})\bigr)=\frac1a+\frac1{a^3}+O(a^{-5})$, using $\varepsilon=\frac1a+O(a^{-3})$, itself obtained from $\varepsilon=\arctan\frac1{a-\varepsilon}=\frac1a+O(a^{-3})$. Finally
$\varepsilon=\arctan y=y-\frac{y^3}3+O(y^5)$ with $y=\frac1a+\frac1{a^3}+O(a^{-5})$ gives $\varepsilon=\frac1a+\frac2{3a^3}+O(a^{-5})$.

(c) The expansion is in powers of $\frac1a$, and $a=\frac{3\pi}2=4.71$ already for $k=1$; the next term, computed in the notebook, is $-\frac{13}{15a^5}$, about $-0.0004$, which is the
observed error.

\tools{Tome~0, Chapter~6 (derivatives, Taylor expansions of $\tan$ and $\arctan$).}
\leads These roots are the zeros of $\frac{\sin x-x\cos x}{x^3}=\frac13-\frac{x^2}{30}+\cdots$. A product expansion like Euler's sine product then gives $\sum_k\frac1{x_k^2}=\frac{1/30}{1/3}=\frac1{10}$, which the
notebook confirms. In practice such expansions are excellent starting points for Newton's method.
\end{solution}

\begin{solution}{pb:nbonacci}
\idea Multiplying by $x-1$ makes the equation $x^n(2-x)=1$: the root is $2$ minus something exponentially small, and the equation for that small thing can be iterated.

\writeup (a) $P_n(x)=x^n\bigl(1-\sum_{j=1}^nx^{-j}\bigr)$, and the bracket increases strictly on $\left]0,\infty\right[$ from $-\infty$ to $1$: exactly one positive root $r_n$. $P_n(1)=1-n<0$ and $P_n(2)=2^n-(2^n-1)=1>0$, so
$1<r_n<2$. Since $P_{n+1}(x)=xP_n(x)-1$, $P_{n+1}(r_n)=-1<0$, so $r_{n+1}>r_n$. And $r_2=\frac{1+\sqrt5}2$.

(b) Since $(x-1)P_n(x)=x^{n+1}-2x^n+1$, $r=r_n$ satisfies $r^n(2-r)=1$. Let $\varepsilon=2-r\in\left]0,2-r_2\right]$. Then $\varepsilon=(2-\varepsilon)^{-n}\le r_2^{-n}$, so $n\varepsilon\to0$. Now
$(2-\varepsilon)^{-n}=2^{-n}\exp\bigl(-n\ln(1-\frac\varepsilon2)\bigr)=2^{-n}\bigl(1+\frac{n\varepsilon}2+O(n^2\varepsilon^2)\bigr)$. This first gives $\varepsilon=2^{-n}(1+O(n\varepsilon))=O(2^{-n})$, then
$\varepsilon=2^{-n}+n2^{-n-1}\varepsilon+O(n^22^{-3n})=2^{-n}+n2^{-2n-1}+O(n^22^{-3n})$.

(c) A sequence without $n$ consecutive heads ends with a tail preceded by $0,\dots,n-1$ heads, or is shorter than $n$; this gives the recurrence, whose characteristic polynomial is
$P_n$. With $r=r_n$, $\sum_{i=1}^nr^{-i}=1$. If $a_j\ge r^{j-n}$ for $n$ consecutive $j$, then $a_\ell\ge\sum_{i=1}^nr^{\ell-i-n}=r^{\ell-n}$; the initial values $a_j=2^j$ ($j<n$) satisfy this, so by
induction $a_\ell\ge r^{\ell-n}$ for all $\ell$. In the same way $a_\ell\le c\,r^\ell$ with $c=\max_{j<n}2^jr^{-j}$. So $a_\ell^{1/\ell}\to r_n$. And $r_n$ is close to $2$ because forbidding a run of $n$
heads costs little: among the $2^\ell$ sequences, a fraction about $e^{-\ell2^{-n-1}}$ avoids such runs, consistent with $(\frac{r_n}2)^\ell\approx(1-2^{-n-1})^\ell$.

\tools{Tome~0, Chapter~5 (sequences, recurrences); Tome~3, Chapter~4 (runs of heads).}
\leads For $n=3$, $r_3=1.8393$ is the tribonacci constant. The estimate $(1-2^{-n-1})^\ell$ is the heart of the Erd\H os--R\'enyi law of the longest run (Tome~3): runs of length $n$ appear once
$\ell2^{-n}$ is large, that is $n\approx\log_2\ell$.
\end{solution}

\begin{solution}{pb:sqrtsum}
\idea Proposition~11.7 of Tome~4 is an identity between holomorphic functions, valid far to the left of $\operatorname{Re}s=1$. Read at $s=-\frac12$ it is exactly Euler--Maclaurin for $\sum\sqrt k$, with $\zeta(-\frac12)$
in the role of the constant.

\writeup Tome~4, Proposition~11.7 states that for $N,M\ge1$, $\operatorname{Re}s>1-2M$, $s\ne1$,
\[
\zeta(s)=\sum_{k\le N}k^{-s}+\frac{N^{1-s}}{s-1}-\frac{N^{-s}}2+\sum_{j=1}^M\frac{B_{2j}}{(2j)!}(s)_{2j-1}N^{1-s-2j}-\frac{(s)_{2M}}{(2M)!}\int_N^\infty\frac{\bar B_{2M}(x)}{x^{s+2M}}\dd x,
\]
with $(s)_i=s(s+1)\cdots(s+i-1)$ and $\bar B_{2M}$ bounded. Take $M=2$. For $s=-\frac12$: $\frac{N^{3/2}}{-3/2}=-\frac23N^{3/2}$, $-\frac{N^{1/2}}2$, and the $j=1$ term is $\frac1{12}\cdot(-\frac12)N^{-1/2}=-\frac1{24\sqrt N}$; the $j=2$ term
and the integral are $O(N^{-5/2})$. So $\sum_{k\le N}\sqrt k=\zeta(-\frac12)+\frac23N^{3/2}+\frac12\sqrt N+\frac1{24\sqrt N}+O(N^{-5/2})$. For $s=\frac12$: $\frac{N^{1/2}}{-1/2}=-2\sqrt N$, $-\frac1{2\sqrt N}$, the $j=1$ term is
$\frac1{12}\cdot\frac12N^{-3/2}=\frac1{24N^{3/2}}$, and the rest is $O(N^{-7/2})$. So $\sum_{k\le N}\frac1{\sqrt k}=\zeta(\frac12)+2\sqrt N+\frac1{2\sqrt N}-\frac1{24N^{3/2}}+O(N^{-7/2})$.

\tools{Tome~4, Theorem~11.5 (Euler--Maclaurin) and Proposition~11.7 (continuation of $\zeta$).}
\leads $\zeta(\frac12)<0$ although every term is positive: the constant is what is left after removing $2\sqrt N$, not a sum. By the functional equation (Tome~4, Chapter~10),
$\zeta(-\frac12)=-\frac{\zeta(3/2)}{4\pi}$. Ramanujan called such constants ``the constant of a series'', and Hardy's \emph{Divergent Series} (1949) makes the idea systematic.
\end{solution}

\begin{solution}{pb:hyperfactorial}
\idea $\ln H(N)=\sum k\ln k$ is minus the derivative of $\sum k^{-s}$ at $s=-1$. Differentiating the Euler--Maclaurin identity for $\zeta$ therefore gives the expansion of $\ln H$, with $\zeta'(-1)$ as
the constant.

\writeup Take $M=3$ in the identity above (valid for $\operatorname{Re}s>-5$) and differentiate in $s$; both sides are holomorphic there, and the integral may be differentiated under the
integral sign. At $s=-1$:
\begin{itemize}
\item $\frac{\mathrm d}{\mathrm ds}\sum_{k\le N}k^{-s}=-\sum k^{-s}\ln k$, which is $-\ln H(N)$;
\item $\frac{\mathrm d}{\mathrm ds}\frac{N^{1-s}}{s-1}=-\frac{N^{1-s}\ln N}{s-1}-\frac{N^{1-s}}{(s-1)^2}$, which is $\frac{N^2}2\ln N-\frac{N^2}4$;
\item $\frac{\mathrm d}{\mathrm ds}\bigl(-\frac{N^{-s}}2\bigr)=\frac{N^{-s}\ln N}2$, which is $\frac N2\ln N$;
\item $\frac{\mathrm d}{\mathrm ds}\bigl(\frac{B_2}2sN^{-1-s}\bigr)=\frac1{12}N^{-1-s}(1-s\ln N)$, which is $\frac1{12}(1+\ln N)$;
\item for $j=2,3$ the factor $(s)_{2j-1}$ vanishes at $s=-1$, so only its derivative survives: $(s)_3'(-1)=(-1)\cdot1=-1$, giving $\frac{B_4}{24}(-1)N^{-2}=\frac1{720N^2}$, and the $j=3$ term gives
$O(N^{-4})$;
\item likewise $(s)_6$ vanishes at $-1$, so the derivative of the remainder is
\[
-\frac{(s)_6'(-1)}{6!}\int_N^\infty\bar B_6(x)x^{-5}\dd x=O(N^{-4}).
\]
\end{itemize}
So
\[
\begin{aligned}\zeta'(-1)&=-\ln H(N)+\Bigl(\frac{N^2}2+\frac N2+\frac1{12}\Bigr)\ln N\\&\quad-\frac{N^2}4+\frac1{12}+\frac1{720N^2}+O(N^{-4}),\end{aligned}
\]
which is the claim with $\ln A=\frac1{12}-\zeta'(-1)$.

(b) $\zeta'(-1)=-0.1654211\ldots$, so $\ln A=0.2487545$ and $A=1.2824271$. For $n=10$, $\ln H(10)=\sum_{k\le10}k\ln k=102.0828\ldots$ and the formula without the $\frac1{720n^2}$ term is off by $1.39\cdot10^{-5}$, which is $\frac1{72000}$ to
three digits.

\tools{Tome~4, Proposition~11.7 and Chapter~9 (differentiating holomorphic functions).}
\leads Kinkelin (1860) and Glaisher (1878) found this constant; Glaisher's name stuck. With the functional equation one can show $\zeta'(2)=\frac{\pi^2}6(\gamma+\ln2\pi-12\ln A)$, so $A$ also
measures how $\zeta$ changes at $2$. $A$ is to $\prod k^k$ what $\sqrt{2\pi}$ is to $\prod k$: indeed $\ln\sqrt{2\pi}=-\zeta'(0)$, and the same differentiation at $s=0$ gives Stirling's formula.
\end{solution}

\begin{solution}{pb:diehit}
\idea The generating function is rational. Its only pole on the closed unit disc is at $1$, with residue $\frac1{\text{mean}}$; the other poles are farther out and only produce an
exponentially small error.

\writeup (a) Condition on the first roll $j$: the totals then hit $n$ if and only if the later totals, started afresh, hit $n-j$. So $u_n=\frac16\sum_ju_{n-j}$ for $n\ge1$, and $U=\sum u_nz^n$ satisfies
$U=1+D(z)U$.

(b) $D(1)=1$ and $D'(1)=\frac{1+\dots+6}6=\frac72\ne0$, so $1$ is a simple root. If $|z|\le1$ and $D(z)=1$, then $1=|D(z)|\le\frac16\sum|z|^j\le1$, so $|z|=1$ and all the $z^j$ have the same argument
as their sum, $1$: $z^j=1$ for $j=1,\dots,6$, so $z=1$.

(c) The polynomial $1-D$ factors as $(1-z)q(z)$ with $q(1)=D'(1)=\frac72$, and its other roots have modulus at least some $R>1$. So $V(z)=U(z)-\frac{2/7}{1-z}$ is rational without poles in
$|z|<R$, and Cauchy's estimates (Tome~4, Corollary~9.12) give $[z^n]V=O(r^{-n})$ for each $r<R$. Hence $u_n=\frac27+O(r^{-n})$. The notebook finds $R=1.3694$, so $\rho=0.7302$.

(d) The same argument applies to $D(z)=\sum_{j\le N}p_jz^j$: equality $D(z)=1$ on $|z|\le1$ forces $z^j=1$ for every $j$ with $p_j>0$, hence $z^g=1$ with $g$ the gcd of these $j$ (by B\'ezout, $g$ is an integer
combination of them), so $z=1$ when $g=1$; and $D'(1)=\mu$. If $g=2$, all totals are even, $u_n=0$ for odd $n$, and $z=-1$ is a second pole on the circle: $u_{2n}\to\frac2\mu$ instead.

\tools{Tome~4, Chapter~2 (generating functions) and Corollary~9.12 (Cauchy's estimates); Tome~3, Chapter~7 (first-step decompositions).}
\leads $\frac27$ is the renewal theorem: hitting points are spread at average spacing $\mu$. Erd\H os, Feller and Pollard (1949) proved $u_n\to\frac1\mu$ for any step law on the positive integers
with gcd $1$, finite support or not; without finite support there is no rational function and no exponential rate, and the proof needs a Tauberian idea.
\end{solution}

\begin{solution}{pb:banach}
\idea Finding a box empty means asking for it one time too many. The recursion between consecutive probabilities sums to an equation for the mean, which leaves a central
binomial coefficient.

\writeup (a) Suppose the left box is found empty. This happens when the left pocket is chosen for the $(N+1)$-th time, the right one having been chosen $N-r$ times before: $2N-r+1$
choices, the last one left. This has probability $\binom{2N-r}N2^{-(2N-r+1)}$, and doubling for the right box gives the formula.

(b) $\frac{a_{r+1}}{a_r}=2\binom{2N-r-1}N\binom{2N-r}N^{-1}=\frac{2(N-r)}{2N-r}$, which is the relation (also true at $r=N$ with $a_{N+1}=0$). Write it $2Na_{r+1}-ra_{r+1}=2Na_r-2ra_r$ and sum over
$0\le r\le N$. With $E=\mathbb E[R]=\sum ra_r$ and $\sum a_r=1$: the left side is $2N(1-a_0)-\sum_{r'\ge1}(r'-1)a_{r'}=2N(1-a_0)-E+(1-a_0)$, the right side is $2N-2E$. So $E=(2N+1)a_0-1$, with
$a_0=\binom{2N}N4^{-N}$.

(c) $a_0=\frac1{\sqrt{\pi N}}(1+O(\frac1N))$, so $E=\frac{2N+1}{\sqrt{\pi N}}-1+O(N^{-1/2})=2\sqrt{\frac N\pi}-1+O(N^{-1/2})$. For $N=50$: $E=7.04$, against $2\sqrt{50/\pi}-1=6.98$.

\tools{Tome~3, Chapter~2 (expectations); Tome~4, Proposition~11.12 (central binomial coefficients).}
\leads Feller's textbook made this ``Banach match box problem'' famous. The answer $2\sqrt{N/\pi}$ is the expected distance of a simple random walk from its start after about $2N$ steps,
which is what $R$ really is: the difference between the two boxes when one runs out.
\end{solution}

\begin{solution}{pb:flory}
\idea The first dimer cuts the row into two independent rows: a recursion with a uniform split, which becomes a first-order differential equation for the generating
function, solved by an exponential.

\writeup (a) The pair with the smallest label is filled first, and it is uniform among the $n-1$ pairs. Given that it is $\{i,i+1\}$, the labels of the other pairs are independent and
uniform on $]\ell,1]$ ($\ell$ the smallest label), and the pairs touching $i$ or $i+1$ can never be filled; the process continues independently on the rows $1,\dots,i-1$ and
$i+2,\dots,n$, each with the same rule (rescaling labels does not change their order). So $V_n=\frac1{n-1}\sum_{i=1}^{n-1}(V_{i-1}+V_{n-i-1})=\frac2{n-1}\sum_{j=0}^{n-2}V_j$.

(b) The relation $(n-1)V_n=2\sum_{j\le n-2}V_j$ also holds for $n=0,1$ (both sides vanish). Multiply by $x^n$ and sum: $xV'-V=\frac{2x^2V}{1-x}$, where $V=\sum V_nx^n$. Put $V=xW$:
$x^2W'=\frac{2x^3W}{1-x}$, so $\frac{W'}W=\frac{2x}{1-x}=-2+\frac2{1-x}$ and $W=\frac{ce^{-2x}}{(1-x)^2}$, with $c=W(0)=V_1=1$. The coefficient of $x^{n-1}$ in $\frac{e^{-2x}}{(1-x)^2}=\sum_k\frac{(-2)^kx^k}{k!}\sum_m(m+1)x^m$ is
$\sum_{k\le n-1}\frac{(-2)^k}{k!}(n-k)$.

(c) The full series is $\sum_{k\ge0}(n-k)\frac{(-2)^k}{k!}=ne^{-2}+2e^{-2}$, so $V_n-(n+2)e^{-2}=\sum_{j\ge1}j\frac{(-2)^{n+j}}{(n+j)!}$ (with $k=n+j$; up to sign). Its terms alternate in sign, and their absolute values $t_j$
satisfy $\frac{t_{j+1}}{t_j}=\frac{2(j+1)}{j(n+j+1)}\le\frac4{n+2}\le1$ for $n\ge2$, so the sum is at most $t_1=\frac{2^{n+1}}{(n+1)!}$ in absolute value. For $n=0,1$ check directly ($|0-2e^{-2}|$ and $|1-3e^{-2}|$ are less than $2$).

(d) $V=h(x)(1-x)^{-2}$ with $h(x)=xe^{-2x}$ entire, so Tome~4, Theorem~11.13 with $\alpha=2$ gives $V_n=\frac{h(1)}{\Gamma(2)}n+O(1)=e^{-2}n+O(1)$, as in (c).

\tools{Tome~4, Chapter~2 (generating functions, differential equations) and Theorem~11.13 (transfer); Tome~3, Chapter~5 (uniform labels).}
\leads Paul Flory (1939) computed this fraction for chemical groups reacting in pairs along a polymer, and Page (1959) the distribution of the vacancies. In the continuous
version, cars of length $1$ park at random on a street of length $x$, and R\'enyi (1958) showed that they fill a fraction tending to $0.7476\ldots$; the notebook simulates both.
\end{solution}

\begin{solution}{pb:saw}
\idea A long self-avoiding walk is two shorter ones glued together, but not every gluing is self-avoiding: submultiplicativity, hence a growth rate by Fekete. Simple
families of walks squeeze it.

\writeup (a) Cutting a self-avoiding walk of length $m+n$ after $m$ steps gives a self-avoiding walk of length $m$ and, translated to the origin, one of length $n$; the map is injective.

(b) Fix $k$ and write $n=qk+r$ with $0\le r<k$. Iterating, $a_n\le qa_k+a_r$ (with $a_0$ replaced by $0$ when $r=0$). So $\frac{a_n}n\le\frac{qk}n\cdot\frac{a_k}k+\frac{\max_{r<k}|a_r|}n$, and $\limsup\frac{a_n}n\le\frac{a_k}k$ for every
$k$. Hence $\limsup\frac{a_n}n\le\inf_k\frac{a_k}k\le\liminf\frac{a_n}n$.

(c) Apply (b) to $a_n=\ln c_n$.

(d) After the first step a self-avoiding walk never steps straight back, so $c_n\le4\cdot3^{n-1}$ and $\mu\le3$. For the lower bound, a walk with steps east, north, south and no north--south
reversal is self-avoiding: its $x$-coordinate never decreases, and at a given $x$ it moves monotonically along one vertical segment. Let $e_n$ and $v_n$ count such walks of length
$n$ ending with a horizontal or vertical step. Then $e_{n+1}=e_n+v_n$ and $v_{n+1}=2e_n+v_n$, with $e_1=1$, $v_1=2$. The matrix $\begin{psmallmatrix}1&1\\2&1\end{psmallmatrix}$ has eigenvalues $1\pm\sqrt2$, with eigenvectors $(1,\pm\sqrt2)$, and
$(1,2)=\frac{1+\sqrt2}2(1,\sqrt2)+\frac{1-\sqrt2}2(1,-\sqrt2)$ has a positive component along the first. So $e_n+v_n\ge c(1+\sqrt2)^n$
with $c>0$, and $\mu\ge1+\sqrt2$.

(e) $\mu\le c_{14}^{1/14}=2374444^{1/14}=2.854$.

\tools{Tome~0, Chapter~7 (eigenvalues of $2\times2$ matrices); Tome~0, Chapter~5 (lim sup and lim inf).}
\leads Hammersley and Morton (1954) proved that $\mu$ exists. Its value on $\Z^2$ is unknown; numerical work gives $\mu=2.6381585\ldots$ On the hexagonal lattice, Nienhuis (1982) predicted
$\mu=\sqrt{2+\sqrt2}$, and Duminil-Copin and Smirnov (2012) proved it. The ratios $\frac{c_{n+1}}{c_n}$ in Figure~\ref{fig:saw} oscillate with the parity of $n$ but approach $\mu$ much faster than $c_n^{1/n}$.
\end{solution}

\begin{solution}{pb:lcs}
\idea Common subsequences of two halves glue into a common subsequence of the whole, so the mean is superadditive and has a linear growth rate. The first moment
method caps it.

\writeup (a) If $A=A'A''$ and $B=B'B''$ with $|A'|=|B'|=m$, a common subsequence of $(A',B')$ followed by one of $(A'',B'')$ is common to $(A,B)$. The pairs $(A',B')$ and $(A'',B'')$ have the laws of
random pairs of lengths $m$ and $n$, so $\mathbb E[L_{m+n}]\ge\mathbb E[L_m]+\mathbb E[L_n]$. Since $\mathbb E[L_n]\le n$, Fekete's lemma (Problem~\ref{pb:saw}) applied to $-\mathbb E[L_n]$ gives convergence of $\frac{\mathbb E[L_n]}n$ to
$\sup_n\frac{\mathbb E[L_n]}n\le1$.

(b) $\gamma_2\ge\frac{\mathbb E[L_8]}8=\frac{44803}{65536}=0.6836$.

(c) $L_n\ge k$ means that for some index sets $I,J\subset\{1,\dots,n\}$ of size $k$ the subwords $A_I$ and $B_J$ are equal; for fixed $I,J$ they are independent uniform words of length $k$, equal with
probability $2^{-k}$. The union bound gives $P(L_n\ge k)\le\binom nk^22^{-k}$. From $1\ge\binom nkx^k(1-x)^{n-k}$ at $x=\frac kn$ we get $\binom nk\le e^{nH(k/n)}$. The function $\lambda(\kappa)=2H(\kappa)-\kappa\ln2$ is
concave, positive at $\frac12$, equal to $-\ln2$ at $1$, and decreasing on $[\frac12,1]$ (its derivative is $2\ln\frac{1-\kappa}\kappa-\ln2<0$), so it has one root $\kappa^*$ there. For $\kappa>\kappa^*$ and $k=\lceil\kappa n\rceil$,
$P(L_n\ge k)\le e^{n\lambda(\kappa)}$, and $\mathbb E[L_n]\le k+nP(L_n\ge k)\le\kappa n+1+ne^{n\lambda(\kappa)}$. Dividing by $n$, $\gamma_2\le\kappa$; so $\gamma_2\le\kappa^*=0.9051$.

\tools{Tome~3, Chapter~2 (union bound, expectations) and the binomial theorem; Problem~\ref{pb:saw} (Fekete).}
\leads Chv\'atal and Sankoff (1975) proved that $\gamma_2$ exists. Its value is still unknown: Lueker (2009) proved $0.788071\le\gamma_2\le0.826280$, and the notebook's simulations suggest about
$0.81$. Longest common subsequences compare files (the \texttt{diff} utility) and DNA sequences, where one wants to know how long a common subsequence two unrelated
sequences have by chance.
\end{solution}

\section*{Map of the chapter}

\begin{center}
\small
\adjustbox{max width=\linewidth}{\begin{tabular}{lll}
Problem & Topic & Tomes used\\\hline
\ref{pb:powersum} & $\sum k^n$, a discrete Laplace sum & 0, 4\\
\ref{pb:involasym} & involutions by the saddle point & 0, 4\\
\ref{pb:bellasym} & the size of the Bell numbers & 0, 4\\
\ref{pb:distinctcycles} & cycles of distinct lengths, $e^{-\gamma}$ & 0, 4\\
\ref{pb:zigzag} & alternating permutations, $\sec+\tan$ & 4\\
\ref{pb:binompowers} & $\sum\binom nk^s$ & 3, 4\\
\ref{pb:poissonise} & a triple birthday & 3, 4\\
\ref{pb:ramanujan} & Ramanujan's $Q$ and $\theta_n$ & 0, 3, 4\\
\ref{pb:dixie} & collecting coupons $m$ times & 0, 3\\
\ref{pb:mappings} & random maps & 3, 4\\
\ref{pb:parking} & parking functions, Abel's identity & 0, 3\\
\ref{pb:xplusinv} & $x_{n+1}=x_n+\frac1{x_n}$ & 0, 4\\
\ref{pb:sinesecond} & the sine iteration, second term & 0\\
\ref{pb:tanroots} & the roots of $\tan x=x$ & 0\\
\ref{pb:nbonacci} & $n$-step Fibonacci roots & 0, 3\\
\ref{pb:sqrtsum} & $\sum\sqrt k$ and $\zeta(-\frac12)$ & 4\\
\ref{pb:hyperfactorial} & $\prod k^k$ and Glaisher's constant & 4\\
\ref{pb:diehit} & hitting a total with a die & 3, 4\\
\ref{pb:banach} & Banach's matchboxes & 3, 4\\
\ref{pb:flory} & random dimers on a line & 3, 4\\
\ref{pb:saw} & self-avoiding walks, Fekete's lemma & 0\\
\ref{pb:lcs} & common subsequences of random words & 3\\
\end{tabular}}
\end{center}

\section*{Sources and further reading}

\begin{itemize}
\item G.~H.~Hardy, \emph{Ramanujan: Twelve Lectures on Subjects Suggested by His Life and Work}, Cambridge University Press, 1940, Lecture~I (the epigraph); the lecture first
appeared as The Indian mathematician Ramanujan, \emph{American Mathematical Monthly} 44 (1937), 137--155.
\item S.~Ramanujan, Question~294, \emph{Journal of the Indian Mathematical Society} 3 (1911), 128; G.~Szeg\H o, \"Uber einige von S.~Ramanujan gestellte Aufgaben, \emph{Journal of the
London Mathematical Society} 3 (1928), 225--232; G.~N.~Watson, Theorems stated by Ramanujan (V): approximations connected with $e^x$, \emph{Proceedings of the London
Mathematical Society} (2) 29 (1929), 293--308; P.~Flajolet, P.~J.~Grabner, P.~Kirschenhofer and H.~Prodinger, On Ramanujan's $Q$-function, \emph{Journal of Computational and Applied
Mathematics} 58 (1995), 103--116.
\item D.~E.~Knuth, Linear probing and graphs, \emph{Algorithmica} 22 (1998), 561--568; D.~E.~Knuth, \emph{The Art of Computer Programming}, Volume~3, \emph{Sorting and Searching}, Addison-Wesley,
2nd edition 1998, Section~6.4; P.~Flajolet, P.~Poblete and A.~Viola, On the analysis of linear probing hashing, \emph{Algorithmica} 22 (1998), 490--515.
\item N.~G.~de Bruijn, \emph{Asymptotic Methods in Analysis}, North-Holland, 1958; P.~Flajolet and R.~Sedgewick, \emph{Analytic Combinatorics}, Cambridge University Press, 2009.
\item W.~K.~Hayman, A generalisation of Stirling's formula, \emph{Journal f\"ur die reine und angewandte Mathematik} 196 (1956), 67--95; S.~Chowla, I.~N.~Herstein and W.~K.~Moore, On recursions
connected with symmetric groups I, \emph{Canadian Journal of Mathematics} 3 (1951), 328--334; L.~Moser and M.~Wyman, On solutions of $x^d=1$ in symmetric groups, \emph{Canadian Journal of
Mathematics} 7 (1955), 159--168; L.~Moser and M.~Wyman, An asymptotic formula for the Bell numbers, \emph{Transactions of the Royal Society of Canada} III 49 (1955), 49--54.
\item D.~Andr\'e, D\'eveloppements de $\sec x$ et de $\tan x$, \emph{Comptes Rendus de l'Acad\'emie des Sciences} 88 (1879), 965--967.
\item J.~Franel, questions in \emph{L'Interm\'ediaire des Math\'ematiciens} 1 (1894) and 2 (1895).
\item M.~S.~Klamkin and D.~J.~Newman, Extensions of the birthday surprise, \emph{Journal of Combinatorial Theory} 3 (1967), 279--282.
\item D.~J.~Newman and L.~Shepp, The double dixie cup problem, \emph{American Mathematical Monthly} 67 (1960), 58--61; P.~Erd\H os and A.~R\'enyi, On a classical problem of probability
theory, \emph{Magyar Tudom\'anyos Akad\'emia Matematikai Kutat\'o Int\'ezet\'enek K\"ozlem\'enyei} 6 (1961), 215--220.
\item M.~D.~Kruskal, The expected number of components under a random mapping function, \emph{American Mathematical Monthly} 61 (1954), 392--397; B.~Harris, Probability distributions
related to random mappings, \emph{Annals of Mathematical Statistics} 31 (1960), 1045--1062; J.~M.~Pollard, A Monte Carlo method for factorization, \emph{BIT} 15 (1975), 331--334.
\item A.~G.~Konheim and B.~Weiss, An occupancy discipline and applications, \emph{SIAM Journal on Applied Mathematics} 14 (1966), 1266--1274; J.~Riordan, Ballots and trees, \emph{Journal of
Combinatorial Theory} 6 (1969), 408--411.
\item H.~Kinkelin, \"Uber eine mit der Gammafunction verwandte Transcendente und deren Anwendung auf die Integralrechnung, \emph{Journal f\"ur die reine und angewandte Mathematik} 57
(1860), 122--138; J.~W.~L.~Glaisher, On the product $1^1\cdot2^2\cdot3^3\cdots n^n$, \emph{Messenger of Mathematics} 7 (1878), 43--47; G.~H.~Hardy, \emph{Divergent Series}, Oxford University Press,
1949.
\item P.~Erd\H os, W.~Feller and H.~Pollard, A property of power series with positive coefficients, \emph{Bulletin of the AMS} 55 (1949), 201--204; W.~Feller, \emph{An Introduction to
Probability Theory and Its Applications}, Volume~I, Wiley, 1950 (3rd edition 1968).
\item P.~J.~Flory, Intramolecular reaction between neighboring substituents of vinyl polymers, \emph{Journal of the American Chemical Society} 61 (1939), 1518--1521; E.~S.~Page, The
distribution of vacancies on a line, \emph{Journal of the Royal Statistical Society B} 21 (1959), 364--374; A.~R\'enyi, On a one-dimensional problem concerning random space filling,
\emph{Publications of the Mathematical Institute of the Hungarian Academy of Sciences} 3 (1958), 109--127.
\item J.~M.~Hammersley and K.~W.~Morton, Poor man's Monte Carlo, \emph{Journal of the Royal Statistical Society B} 16 (1954), 23--38; B.~Nienhuis, Exact critical point and critical
exponents of $O(n)$ models in two dimensions, \emph{Physical Review Letters} 49 (1982), 1062--1065; H.~Duminil-Copin and S.~Smirnov, The connective constant of the honeycomb lattice
equals $\sqrt{2+\sqrt2}$, \emph{Annals of Mathematics} 175 (2012), 1653--1665; J.~L.~Jacobsen, C.~R.~Scullard and A.~J.~Guttmann, On the growth constant for square-lattice self-avoiding
walks, \emph{Journal of Physics A} 49 (2016), 494004.
\item V.~Chv\'atal and D.~Sankoff, Longest common subsequences of two random sequences, \emph{Journal of Applied Probability} 12 (1975), 306--315; G.~S.~Lueker, Improved bounds on the
average length of longest common subsequences, \emph{Journal of the ACM} 56 (2009), article 17.
\end{itemize}
